\begin{proof}[Proof of \cref{obj:prop:identification}]
Fix \(x\in\mathcal X\). Write
\[
p'_x=P'(X=x),
\qquad
\mu'_x(a)=\frac{E_{P'}[\mathbf1\{X=x\}Y(a)]}{p'_x},
\qquad a\in[0,1].
\]
By \cref{obj:ass:stratum-overlap}, both \(p_x\) and \(p'_x\) are strictly positive.

\begin{enumerate}
\item Define the centered observed dose by
\[
V=W-a_0=A-a_0+\sigma Z.
\]
For every Borel set \(B\subseteq\mathbb R\), define
\[
\begin{aligned}
\nu_{0x}(B)&=P(A-a_0\in B\mid X=x),
&
\nu'_{0x}(B)&=P'(A-a_0\in B\mid X=x),\\
M_{0x}(B)&=\frac{P(X=x,V\in B)}{p_x},
&
M'_{0x}(B)&=\frac{P'(X=x,V\in B)}{p'_x}.
\end{aligned}
\]
Equality of the observed laws equates their normalized restrictions to each stratum, and therefore
\[
M_{0x}=M'_{0x}.
\]

Let the Gaussian noise measure be
\[
\Phi_\sigma=\mathcal L(\sigma Z),
\qquad
\Phi_0=\delta_0,
\qquad
\delta_0(B)=\mathbf1\{0\in B\}.
\]
By \cref{obj:ass:gaussian-channel}, this is the centered Gaussian probability measure with variance \(\sigma^2\) under either structural law. For Borel sets \(B_Z,B_t\subseteq\mathbb R\), \cref{obj:ass:error-independence} gives
\[
\begin{aligned}
P(Z\in B_Z,A-a_0\in B_t\mid X=x)
&=\frac{P(Z\in B_Z)\,P(X=x,A-a_0\in B_t)}{p_x}\\
&=P(Z\in B_Z)\,\nu_{0x}(B_t).
\end{aligned}
\]
Scaling the error coordinate consequently yields
\[
\mathcal L_P(A-a_0,\sigma Z\mid X=x)
=\nu_{0x}\otimes\Phi_\sigma.
\]
The same factorization holds under \(P'\). Pushing these product measures forward under addition gives
\[
\begin{aligned}
M_{0x}(B)
&=\int_{\mathbb R}\int_{\mathbb R}
\mathbf1\{t+z\in B\}\,\Phi_\sigma(dz)\,\nu_{0x}(dt)
=(\nu_{0x}*\Phi_\sigma)(B),\\
M'_{0x}(B)&=(\nu'_{0x}*\Phi_\sigma)(B).
\end{aligned}
\]

The density representation in \cref{obj:ass:weak-design}, together with the support condition in \cref{obj:ass:latent-support} and \(a_0=1/2\), gives
\[
\nu_{0x}(B)=\int_{B\cap[-1/2,1/2]}g_x(t)\,dt.
\]
Both latent laws are probability measures, and the same support conclusion holds under \(P'\):
\[
\nu_{0x}(\mathbb R)=\nu'_{0x}(\mathbb R)=1,
\qquad
\nu_{0x}\bigl(\mathbb R\setminus[-1/2,1/2]\bigr)
=\nu'_{0x}\bigl(\mathbb R\setminus[-1/2,1/2]\bigr)=0.
\]
Thus
\[
\nu_{0x}*\Phi_\sigma=\nu'_{0x}*\Phi_\sigma
\quad\Longrightarrow\quad
\nu_{0x}=\nu'_{0x},
\]
by \cref{obj:lem:compact-gaussian-convolution-injective}, taking its two auxiliary measures to be zero. Its finiteness, support, and error-scale hypotheses are satisfied, including at \(\sigma=0\).
% lean: CausalSmith.Stat.NoisydoseWeakdesignTransition.latentDoseLaw_eq_of_obsLaw_eq

\item Since \(\beta>0\), the Hölder means are continuous on \([0,1]\). Define, for every Borel \(B\subseteq\mathbb R\),
\[
\begin{aligned}
\nu_{1x}(B)
&=\int_{B\cap[-1/2,1/2]}
\max\{\mu_x(t+a_0),0\}\,\nu_{0x}(dt),\\
\nu'_{1x}(B)
&=\int_{B\cap[-1/2,1/2]}
\max\{\mu'_x(t+a_0),0\}\,\nu'_{0x}(dt).
\end{aligned}
\]
The support of the unmarked laws makes these the marked latent measures in the statement. On this support, \cref{obj:ass:mean-range} makes the positive parts equal to the means themselves. Continuity and boundedness ensure integrability, and hence
\[
\begin{aligned}
\nu_{1x}(\mathbb R)
&=\int_{[-1/2,1/2]}\mu_x(t+a_0)\,\nu_{0x}(dt)\le1,
&
\nu'_{1x}(\mathbb R)&\le1,\\
\nu_{1x}\bigl(\mathbb R\setminus[-1/2,1/2]\bigr)
&=0,
&
\nu'_{1x}\bigl(\mathbb R\setminus[-1/2,1/2]\bigr)&=0.
\end{aligned}
\]

By \cref{obj:lem:realized-dose-mean},
\[
0\le Y\le1
\qquad P\text{-almost surely and }P'\text{-almost surely}.
\]
We may therefore define the finite nonnegative observed marked measures by
\[
M_{1x}(B)=E_P[\mathbf1\{V\in B\}Y\mid X=x],
\qquad
M'_{1x}(B)=E_{P'}[\mathbf1\{V\in B\}Y\mid X=x].
\]
They are determined by the normalized observed laws, so
\[
M_{1x}=M'_{1x}.
\]

For a fixed Borel \(B\), use the bounded measurable test
\[
(a,x',z)\longmapsto
\mathbf1\{x'=x,\ a+\sigma z-a_0\in B\},
\qquad
(a,x',z)\in\mathbb R\times\mathcal X\times\mathbb R.
\]
The bounded-test identity in \cref{obj:lem:realized-dose-mean} gives
\[
\begin{aligned}
M_{1x}(B)
&=\frac{E_P[\mathbf1\{X=x,V\in B\}Y]}{p_x}\\
&=\frac{E_P[\mathbf1\{X=x,V\in B\}\mu_X(A)]}{p_x}\\
&=E_P[\mathbf1\{V\in B\}\mu_x(A)\mid X=x].
\end{aligned}
\]
Using the conditional product law established above, and the supported latent mean, we obtain
\[
\begin{aligned}
M_{1x}(B)
&=\int_{[-1/2,1/2]}\int_{\mathbb R}
\mathbf1\{t+z\in B\}\mu_x(t+a_0)\,
\Phi_\sigma(dz)\,\nu_{0x}(dt)\\
&=(\nu_{1x}*\Phi_\sigma)(B).
\end{aligned}
\]
The same calculation under \(P'\) yields
\[
M'_{1x}=\nu'_{1x}*\Phi_\sigma.
\]
Consequently,
\[
\nu_{1x}*\Phi_\sigma=\nu'_{1x}*\Phi_\sigma
\quad\Longrightarrow\quad
\nu_{1x}=\nu'_{1x},
\]
again by \cref{obj:lem:compact-gaussian-convolution-injective} with the two auxiliary measures zero. The displayed finiteness and support properties discharge its measure hypotheses.
% lean: CausalSmith.Stat.NoisydoseWeakdesignTransition.latentMeanMeasure_eq_of_obsLaw_eq

\item For every integer \(j\ge0\), put
\[
\varepsilon_j=\frac1{j+3},
\qquad
B_j=(-\varepsilon_j,\varepsilon_j).
\]
Then
\[
0<\varepsilon_j\le\frac12,
\qquad
B_j\subseteq[-1/2,1/2].
\]
The lower envelope in \cref{obj:ass:weak-design} implies
\[
g_x(t)\ge c_{\mathrm{lo}}|t|^\kappa>0
\qquad
\text{for Lebesgue-almost every }t\in B_j\setminus\{0\}.
\]
The singleton \(\{0\}\) has Lebesgue measure zero, whereas
\[
\int_{B_j}1\,dt=2\varepsilon_j>0.
\]
If the integral of the nonnegative density over \(B_j\) were zero, the density would vanish almost everywhere there, contradicting this positivity. Thus
\[
\nu_{0x}(B_j)=\int_{B_j}g_x(t)\,dt>0.
\]
This proves every asserted shrinking-neighborhood mass is strictly positive.
% lean: CausalSmith.Stat.NoisydoseWeakdesignTransition.latentDoseLaw_local_mass_pos

\item For \(t\in B_j\), both \(t+a_0\) and \(a_0\) belong to \([0,1]\). The radius-one Hölder bound in \cref{obj:ass:holder-mean} gives
\[
|\mu_x(t+a_0)-\mu_x(a_0)|
\le |t|^\beta
\le\varepsilon_j^\beta.
\]
Integrating this pointwise bound against the latent probability law yields
\[
\left|
\int_{B_j}\bigl(\mu_x(t+a_0)-\mu_x(a_0)\bigr)\,\nu_{0x}(dt)
\right|
\le\varepsilon_j^\beta\,\nu_{0x}(B_j).
\]
Since the marked mass equals the integral of the mean, and the denominator is positive,
\[
\begin{aligned}
\left|
\frac{\nu_{1x}(B_j)}{\nu_{0x}(B_j)}-\mu_x(a_0)
\right|
&=
\frac{
\left|\displaystyle\int_{B_j}
\bigl(\mu_x(t+a_0)-\mu_x(a_0)\bigr)\,\nu_{0x}(dt)\right|
}{\nu_{0x}(B_j)}\\
&\le\varepsilon_j^\beta.
\end{aligned}
\]
% lean: CausalSmith.Stat.NoisydoseWeakdesignTransition.latentLocalRatio_error
Finally,
\[
\varepsilon_j\longrightarrow0,
\qquad
\varepsilon_j^\beta\longrightarrow0,
\]
because \(\beta>0\). Squeezing the nonnegative absolute error therefore proves
\[
\lim_{j\to\infty}\frac{\nu_{1x}(B_j)}{\nu_{0x}(B_j)}
=\mu_x(a_0).
\]
% lean: CausalSmith.Stat.NoisydoseWeakdesignTransition.latentLocalRatio_tendsto

\item Projecting the common observed law onto its stratum coordinate gives
\[
p_x=P(X=x)=P'(X=x)=p'_x.
\]
% lean: CausalSmith.Stat.NoisydoseWeakdesignTransition.strataProb_eq_of_obsLaw_eq

\item The equalities of the unmarked and marked latent measures imply, by uniqueness of densities with respect to their common unmarked probability measure,
\[
\max\{\mu_x(t+a_0),0\}
=\max\{\mu'_x(t+a_0),0\}
\qquad
\text{for }\nu_{0x}\text{-almost every }t\in[-1/2,1/2].
\]
The mean-range bounds remove the positive parts, giving
\[
\mu_x(t+a_0)=\mu'_x(t+a_0)
\qquad
\text{for }\nu_{0x}\text{-almost every }t\in[-1/2,1/2].
\]

The lower design envelope makes \(g_x(t)>0\) for Lebesgue-almost every \(t\) in the full latent interval: the possible exceptional point \(t=0\) is null. Thus, for every Borel \(B\subseteq\mathbb R\),
\[
\nu_{0x}(B)=\int_{B\cap[-1/2,1/2]}g_x(t)\,dt=0
\quad\Longrightarrow\quad
\int_{B\cap[-1/2,1/2]}1\,dt=0.
\]
Indeed, a zero integral of the nonnegative density forces it to vanish almost everywhere on that set, and its almost-everywhere positivity forces the set to be Lebesgue-null. The mean equality consequently transfers to Lebesgue-almost every centered dose:
\[
\mu_x(t+a_0)=\mu'_x(t+a_0)
\qquad
\text{for Lebesgue-almost every }t\in[-1/2,1/2].
\]

For \(t,t'\in[-1/2,1/2]\), the two Hölder bounds give
\[
\begin{aligned}
|\mu_x(t+a_0)-\mu_x(t'+a_0)|&\le|t-t'|^\beta,\\
|\mu'_x(t+a_0)-\mu'_x(t'+a_0)|&\le|t-t'|^\beta.
\end{aligned}
\]
Since \(\beta>0\), both centered mean functions are continuous on the closed interval. If they differed at any point, continuity would preserve their difference on a relative neighborhood of that point. Every such neighborhood has positive Lebesgue measure, including at either endpoint, contradicting their almost-everywhere agreement. Hence
\[
\mu_x(t+a_0)=\mu'_x(t+a_0)
\qquad
\text{for every }t\in[-1/2,1/2].
\]
For every \(a\in[0,1]\),
\[
a-a_0\in[-1/2,1/2],
\qquad
(a-a_0)+a_0=a.
\]
Substitution therefore proves
\[
\mu_x(a)=\mu'_x(a)
\qquad
\text{for every }a\in[0,1].
\]
Since \(x\) was arbitrary, all the preceding conclusions hold in both strata.
% lean: CausalSmith.Stat.NoisydoseWeakdesignTransition.condPotMean_eqOn_of_latent_measures_eq

\item Bounded potential outcomes ensure integrability at \(a_0\). The two strata partition the structural space, so pointwise
\[
Y(a_0)=\sum_{x\in\mathcal X}\mathbf1\{X=x\}Y(a_0).
\]
Taking expectations and using \cref{obj:def:conditional-potential-mean}, with the positive stratum probabilities, gives
\[
\begin{aligned}
\theta(P)
&=E_P[Y(a_0)]\\
&=\sum_{x\in\mathcal X}E_P[\mathbf1\{X=x\}Y(a_0)]\\
&=\sum_{x\in\mathcal X}p_x\mu_x(a_0).
\end{aligned}
\]
The same decomposition applies under \(P'\). Since \(a_0=1/2\in[0,1]\), the identified stratum probabilities and mean functions yield
\[
\theta(P')
=\sum_{x\in\mathcal X}p'_x\mu'_x(a_0)
=\sum_{x\in\mathcal X}p_x\mu_x(a_0)
=\theta(P).
\]
% lean: CausalSmith.Stat.NoisydoseWeakdesignTransition.causalTarget_eq_stratum_sum
% lean: CausalSmith.Stat.NoisydoseWeakdesignTransition.identification
\end{enumerate}
\end{proof}

\begin{proof}[Proof of \cref{obj:thm:uniform-frontier}]
Fix public constants \(K=(p_{\min},c_{\mathrm{lo}},c_{\mathrm{hi}})\), \(\beta\in(0,1]\), \(\kappa\in[0,2]\), and \(\alpha\in(0,1/2)\), under the envelope calibration in the statement. Throughout, use the rates and comparison scales defined there. In particular,
\[
d=2\beta+\kappa+1>1>0.
\]

\begin{enumerate}
\item \textbf{Upper risk bound.}
By \cref{obj:thm:observable-upper}, there are \(C_R>0\) and \(n_R\in\mathbb N\), independent of the public path space and of \(\sigma\), such that
\[
\sup_{P\in\mathcal P_{K,\beta,\kappa,\sigma}}
E_{Q_P^n}|\widehat\theta-\theta(P)|
\le C_R\rho_{n,\sigma}
\qquad
(n\ge n_R,\ 0\le\sigma\le1/4).
\]
The measurable estimator in \cref{obj:def:total-estimator} is admissible in \cref{obj:def:minimax-risk}. Viewed as a randomized estimator, it ignores the independent uniform seed, so its risk under the product experiment equals its observed-sample risk. Consequently,
\[
R_n(\sigma)
\le
\sup_{P\in\mathcal P_{K,\beta,\kappa,\sigma}}
E_{Q_P^n}|\widehat\theta-\theta(P)|
\le C_R\rho_{n,\sigma}.
\]
% lean: CausalSmith.Stat.NoisydoseWeakdesignTransition.minimaxRisk_frontier_upper

\item \textbf{Upper length bound.}
By \cref{obj:thm:finite-honesty}, the measurable connected interval in \cref{obj:def:honest-interval} satisfies
\[
Q_P^n\{\theta(P)\in C_{n,\sigma}\}\ge1-\alpha
\]
for every model law, every \(n\), and every allowed \(\sigma\). The same result supplies \(C_H>0\) and \(n_H\in\mathbb N\), uniform over the path spaces and error scales, such that
\[
\sup_{P\in\mathcal P_{K,\beta,\kappa,\sigma}}
E_{Q_P^n}\operatorname{length}(C_{n,\sigma})
\le C_H\rho_{n,\sigma}
\qquad(n\ge n_H).
\]
Thus this procedure belongs to the collection defining \cref{obj:def:honest-length}, and
\[
H_n(\sigma)
\le
\sup_{P\in\mathcal P_{K,\beta,\kappa,\sigma}}
E_{Q_P^n}\operatorname{length}(C_{n,\sigma})
\le C_H\rho_{n,\sigma}.
\]
% lean: CausalSmith.Stat.NoisydoseWeakdesignTransition.honestLength_frontier_upper

\item \textbf{A common lower bound.}
Use the published Jacobi localization bound
\citep[Theorem 2.4, equation (2.14)]{PetrushevXu2005},
the local kernel Lipschitz bound
\citep[Theorem 2.2, equation (2.4)]{https://people.math.sc.edu/pencho/Publications/kpx-06-28-06-web.pdf},
Jacobi orthogonality and the positive squared-norm formula
\citep[Table 18.3.1, Jacobi row and caption]{https://dlmf.nist.gov/18.3},
and the positive-real-axis gamma-ratio asymptotic
\citep[equation (5.11.12)]{https://dlmf.nist.gov/5.11.E12}
as the analytic inputs to \cref{obj:thm:observed-lower}.
The packet parameters are admissible: for \(\kappa>0\), they are
\(r=4\) and \(b=(\kappa-1)/2>-1/2\), while the \(\kappa=0\) construction uses the symmetric parameters \((4,4)\).

Set
\[
\tau:=\min\left\{\frac15,\frac{1-2\alpha}{2}\right\}.
\]
Since \(\alpha<1/2\), this choice gives
\[
0<\tau,\qquad \tau\le\frac15,\qquad
\tau\le1-2\alpha-\tau.
\]
That result supplies \(c_{\mathrm w}>0\) and \(n_{\mathrm w}\in\mathbb N\), uniform over the path spaces, such that for every \(n\ge n_{\mathrm w}\) and every allowed \(\sigma\), there are model laws \(P_+,P_-\) with
\[
\theta_\pm:=\theta(P_\pm),\qquad
\delta_{n,\sigma}:=c_{\mathrm w}\rho_{n,\sigma},
\]
\[
|\theta_+-\theta_-|\ge\delta_{n,\sigma},
\qquad
\operatorname{TV}(Q_{P_+}^n,Q_{P_-}^n)\le\tau.
\]
For \(n\ge1\), \(L_n\ge1\); each branch defining \(b_{n,\sigma}\) is nonnegative when \(\sigma\ge0\). Hence
\[
\rho_{n,\sigma}\ge0,\qquad \delta_{n,\sigma}\ge0.
\]

To handle the permitted randomization, define the seed law and product laws by
\[
\nu_{\mathrm{seed}}:=\operatorname{Unif}[0,1],
\qquad
Q_\pm^{n,\mathrm{rand}}:=Q_{P_\pm}^n\otimes\nu_{\mathrm{seed}}.
\]
For any measurable event \(\mathcal V\) in the sample--seed product space, define its section probability by
\[
f_{\mathcal V}(z)
:=
\nu_{\mathrm{seed}}\{u_{\mathrm{seed}}\in[0,1]:
(z,u_{\mathrm{seed}})\in\mathcal V\}.
\]
This is measurable and lies in \([0,1]\). Product integration and the layer-cake identity therefore give
\[
\begin{aligned}
\bigl|Q_+^{n,\mathrm{rand}}(\mathcal V)
      -Q_-^{n,\mathrm{rand}}(\mathcal V)\bigr|
&=
\left|\int f_{\mathcal V}\,dQ_{P_+}^n
      -\int f_{\mathcal V}\,dQ_{P_-}^n\right|\\
&=
\left|\int_0^1
\left[
Q_{P_+}^n\{f_{\mathcal V}\ge v\}
-
Q_{P_-}^n\{f_{\mathcal V}\ge v\}
\right]\,dv\right|\\
&\le \operatorname{TV}(Q_{P_+}^n,Q_{P_-}^n).
\end{aligned}
\]
Taking the supremum over such events yields
\[
\operatorname{TV}(Q_+^{n,\mathrm{rand}},Q_-^{n,\mathrm{rand}})
\le\tau\le\frac15.
\]

For any measurable real-valued estimator \(T\) on this product experiment, define
\[
\mathcal A_\pm
:=
\left\{(z,u_{\mathrm{seed}}):
|T(z,u_{\mathrm{seed}})-\theta_\pm|
\ge\frac{\delta_{n,\sigma}}2\right\}.
\]
The separation and triangle inequality imply
\(\mathcal A_+^c\subseteq\mathcal A_-\). Thus the total-variation bound gives
\[
\begin{aligned}
Q_+^{n,\mathrm{rand}}(\mathcal A_+)
+Q_-^{n,\mathrm{rand}}(\mathcal A_-)
&\ge
Q_+^{n,\mathrm{rand}}(\mathcal A_+)
+Q_-^{n,\mathrm{rand}}(\mathcal A_+^c)\\
&\ge
1-\operatorname{TV}(Q_+^{n,\mathrm{rand}},Q_-^{n,\mathrm{rand}})
\ge\frac45,
\end{aligned}
\]
and consequently
\[
\max\left\{
Q_+^{n,\mathrm{rand}}(\mathcal A_+),
Q_-^{n,\mathrm{rand}}(\mathcal A_-)
\right\}\ge\frac25.
\]
Pointwise domination by the threshold times its indicator now implies
\[
\begin{aligned}
\max_{\pm}E_{Q_\pm^{n,\mathrm{rand}}}|T-\theta_\pm|
&\ge
\frac{\delta_{n,\sigma}}2
\max_{\pm}Q_\pm^{n,\mathrm{rand}}(\mathcal A_\pm)\\
&\ge\frac{\delta_{n,\sigma}}5.
\end{aligned}
\]
These are inequalities between nonnegative integrals, so they also apply to infinite risks. Since both witnesses belong to the model class, taking the infimum over estimators proves
\[
R_n(\sigma)\ge\frac{\delta_{n,\sigma}}5
=\frac{c_{\mathrm w}}5\rho_{n,\sigma}.
\]

Next, let \(C^\circ\) be any measurable connected-interval procedure honest over the model class, and define its coverage events by
\[
\mathcal G_\pm:=\{z:\theta_\pm\in C^\circ(z)\}.
\]
Honesty and total variation imply
\[
Q_{P_+}^n(\mathcal G_+)\ge1-\alpha,
\qquad
Q_{P_+}^n(\mathcal G_-)
\ge Q_{P_-}^n(\mathcal G_-)-\tau
\ge1-\alpha-\tau.
\]
The complementary-event union bound therefore yields
\[
Q_{P_+}^n(\mathcal G_+\cap\mathcal G_-)
\ge1-2\alpha-\tau.
\]
On this intersection, connectedness gives containment of the closed segment joining the targets:
\[
[\min\{\theta_+,\theta_-\},\max\{\theta_+,\theta_-\}]
\subseteq C^\circ(z).
\]
Its Lebesgue length is \(|\theta_+-\theta_-|\), so
\[
\operatorname{length}(C^\circ(z))
\ge
\delta_{n,\sigma}\mathbf1_{\mathcal G_+\cap\mathcal G_-}(z).
\]
Integrating gives
\[
\sup_{P\in\mathcal P_{K,\beta,\kappa,\sigma}}
E_{Q_P^n}\operatorname{length}(C^\circ)
\ge
E_{Q_{P_+}^n}\operatorname{length}(C^\circ)
\ge(1-2\alpha-\tau)\delta_{n,\sigma}.
\]
Taking the infimum over honest procedures proves the corresponding length bound. The argument includes \(\delta_{n,\sigma}=0\).

Define
\[
c_L:=c_{\mathrm w}\tau>0,
\qquad
n_L:=\max\{n_{\mathrm w},1\}.
\]
For \(n\ge n_L\), the inequalities \(\tau\le1/5\) and \(\tau\le1-2\alpha-\tau\), together with \(\delta_{n,\sigma}\ge0\), give
\[
c_L\rho_{n,\sigma}=\tau\delta_{n,\sigma}
\le\frac{\delta_{n,\sigma}}5\le R_n(\sigma),
\]
\[
c_L\rho_{n,\sigma}=\tau\delta_{n,\sigma}
\le(1-2\alpha-\tau)\delta_{n,\sigma}\le H_n(\sigma).
\]
We have obtained, uniformly over the path spaces and allowed error scales,
\[
c_L\rho_{n,\sigma}\le R_n(\sigma),
\qquad
c_L\rho_{n,\sigma}\le H_n(\sigma)
\qquad(n\ge n_L).
\]
% lean: CausalSmith.Stat.NoisydoseWeakdesignTransition.frontier_both_lower

\item \textbf{Comparison at the second threshold.}
Define the positive constant and logarithmic denominator by
\[
\lambda_*:=\log(e+1)>0,
\qquad
\Lambda_n:=\log(e+n b_n^d)\quad(n\ge1).
\]
Because \(L_n\to\infty\) and \(\log L_n/L_n\to0\), choose \(n_E\ge1\) such that, for every \(n\ge n_E\),
\[
L_n\ge4,
\qquad
\frac d2\frac{\log L_n}{L_n}<\frac12.
\]
Since \(b_n=L_n^{-1/2}\) and \(\log n=L_n-1\), logarithmic monotonicity gives
\[
\Lambda_n
\ge\log(n b_n^d)
=L_n-1-\frac d2\log L_n
>\frac{L_n}{2}-1
\ge\frac{L_n}{4}.
\]
Also \(0<b_n\le1\) and \(d\ge0\), whence
\[
e+n b_n^d\le e+n\le2en.
\]
Using \(\log2\le2\), we obtain
\[
\Lambda_n\le\log2+L_n\le4L_n.
\]
Taking square roots and multiplying by \(\sqrt{L_n}\) yields
\[
\frac{\sqrt{L_n}}2\le\sqrt{\Lambda_n}\le2\sqrt{L_n},
\qquad
\frac{L_n}{2}
\le\sqrt{L_n}\sqrt{\Lambda_n}
\le2L_n.
\]
The defining formulas consequently give
\[
F_n(b_n)=\frac1{\sqrt{L_n}\sqrt{\Lambda_n}},
\qquad
\frac1{2L_n}\le F_n(b_n)\le\frac2{L_n}.
\]
Furthermore, \(b_n^2L_n=1\), so
\[
P_n(b_n)=\frac{\lambda_*}{L_n}.
\]
Define
\[
C_E:=\max\left\{\frac2{\lambda_*},\,2\lambda_*\right\}>0.
\]
Then, for \(n\ge n_E\),
\[
F_n(b_n)
\le\frac2{L_n}
\le C_E\frac{\lambda_*}{L_n}
=C_EP_n(b_n),
\]
\[
P_n(b_n)
=\frac{\lambda_*}{L_n}
\le\frac{C_E}{2L_n}
\le C_EF_n(b_n).
\]
% lean: CausalSmith.Stat.NoisydoseWeakdesignTransition.secondElbow_comparison

\item \textbf{Uniform constants and both threshold comparisons.}
Set
\[
c:=c_L,
\qquad
C:=\max\left\{
C_R,C_H,\sqrt{\lambda_*},\frac1{\sqrt{\lambda_*}},
C_E,c_L+1
\right\},
\]
\[
n_0:=\max\{n_R,n_H,n_L,n_E\}.
\]
Thus \(0<c<C\) and \(n_0\ge1\). The second-threshold inequalities just proved remain valid with \(C\) in place of \(C_E\), because the comparison scales are nonnegative.

For \(n\ge n_0\), positivity of \(n\) and \(d\) gives the exact identities
\[
n a_n^d
=n(n^{-1/d})^d
=n n^{-1}=1,
\qquad
F_n(a_n)=\frac{a_n}{\sqrt{\lambda_*}}.
\]
The choices \(C\ge\sqrt{\lambda_*}\) and
\(C\ge1/\sqrt{\lambda_*}\) therefore imply
\[
D_n=a_n\le C F_n(a_n),
\qquad
F_n(a_n)\le C D_n.
\]
For every \(\sigma\in[0,1/4]\), nonnegativity of the branches in
\cref{obj:def:frontier-rate} gives \(\rho_{n,\sigma}\ge0\). Hence enlarging the upper-bound constants preserves their inequalities:
\[
c\rho_{n,\sigma}\le R_n(\sigma)
\le C_R\rho_{n,\sigma}\le C\rho_{n,\sigma},
\]
\[
c\rho_{n,\sigma}\le H_n(\sigma)
\le C_H\rho_{n,\sigma}\le C\rho_{n,\sigma}.
\]
Together with the second-threshold bounds, this establishes the first group of assertions with one \(C\) and one \(n_0\).
% lean: CausalSmith.Stat.NoisydoseWeakdesignTransition.uniform_frontier

\item \textbf{Comparison factors for transition windows.}
Fix a real \(\omega_{\mathrm{win}}\ge1\), and define
\[
M_{\mathrm F,\omega_{\mathrm{win}}}:=\omega_{\mathrm{win}}(1+d\log \omega_{\mathrm{win}}),
\qquad
M_{\mathrm P,\omega_{\mathrm{win}}}:=1+2\log \omega_{\mathrm{win}},
\]
\[
C_{\mathrm a}:=\max\left\{\sqrt{\lambda_*},
\frac1{\sqrt{\lambda_*}}\right\},
\qquad
C_K:=\max\{1,M_{\mathrm F,\omega_{\mathrm{win}}}C_{\mathrm a},
M_{\mathrm F,\omega_{\mathrm{win}}}C_EM_{\mathrm P,\omega_{\mathrm{win}}}\}.
\]
All these factors are positive, and \(C_K\ge1\).

We first derive the scale comparisons used throughout either window. For positive radii \(\varrho_1,\varrho_2\), a nonnegative exponent \(\gamma\), and a nonnegative weight \(\omega_{\log}\), suppose
\[
\varrho_1,\varrho_2>0,\qquad
\gamma,\omega_{\log}\ge0,\qquad
\varrho_1\le \omega_{\mathrm{win}}\varrho_2.
\]
Since \(\omega_{\mathrm{win}}^\gamma\ge1\) and
\(\varrho_1^\gamma\le \omega_{\mathrm{win}}^\gamma\varrho_2^\gamma\),
\[
e+\omega_{\log}\varrho_1^\gamma
\le \omega_{\mathrm{win}}^\gamma(e+\omega_{\log}\varrho_2^\gamma).
\]
Taking logarithms gives
\[
\log(e+\omega_{\log}\varrho_1^\gamma)
\le\gamma\log \omega_{\mathrm{win}}+
\log(e+\omega_{\log}\varrho_2^\gamma).
\]
The logarithm on the right is at least \(1\), while
\(\gamma\log \omega_{\mathrm{win}}\ge0\). Therefore,
\[
\log(e+\omega_{\log}\varrho_1^\gamma)
\le
(1+\gamma\log \omega_{\mathrm{win}})
\log(e+\omega_{\log}\varrho_2^\gamma).
\]
Since \(1+\gamma\log \omega_{\mathrm{win}}\ge1\), this also implies
\[
\log(e+\omega_{\log}\varrho_1^\gamma)
\le
(1+\gamma\log \omega_{\mathrm{win}})^2
\log(e+\omega_{\log}\varrho_2^\gamma),
\]
and taking nonnegative square roots yields
\[
\sqrt{\log(e+\omega_{\log}\varrho_1^\gamma)}
\le
(1+\gamma\log \omega_{\mathrm{win}})
\sqrt{\log(e+\omega_{\log}\varrho_2^\gamma)}.
\]

Now let \(n\ge1\), \(\varrho>0\), and
\(\sigma\in[\varrho/\omega_{\mathrm{win}},\omega_{\mathrm{win}}\varrho]\). Both
\(\sigma\le \omega_{\mathrm{win}}\varrho\) and \(\varrho\le \omega_{\mathrm{win}}\sigma\) hold.
Applying the square-root comparison in both directions with
\(\gamma=d\) and \(\omega_{\log}=n\) gives
\[
\begin{aligned}
F_n(\sigma)
&=
\frac{\sigma}{\varrho}
\frac{\sqrt{\log(e+n\varrho^d)}}
     {\sqrt{\log(e+n\sigma^d)}}F_n(\varrho)
\le M_{\mathrm F,\omega_{\mathrm{win}}}F_n(\varrho),\\
F_n(\varrho)
&=
\frac{\varrho}{\sigma}
\frac{\sqrt{\log(e+n\sigma^d)}}
     {\sqrt{\log(e+n\varrho^d)}}F_n(\sigma)
\le M_{\mathrm F,\omega_{\mathrm{win}}}F_n(\sigma).
\end{aligned}
\]
Applying the logarithmic comparison in both directions with
\(\gamma=2\) and \(\omega_{\log}=L_n>0\), and dividing by \(L_n\), gives
\[
P_n(\sigma)\le M_{\mathrm P,\omega_{\mathrm{win}}}P_n(\varrho),
\qquad
P_n(\varrho)\le M_{\mathrm P,\omega_{\mathrm{win}}}P_n(\sigma).
\]
Thus both scales vary by the stated fixed factors throughout any such radius window.
% lean: CausalSmith.Stat.NoisydoseWeakdesignTransition.fourierScale_window_comparison
% lean: CausalSmith.Stat.NoisydoseWeakdesignTransition.polynomialScale_window_comparison

\item \textbf{One constant and threshold for both windows.}
The positive exponent \(d\) implies
\[
a_n=n^{-1/d}\longrightarrow0,
\qquad
b_n=L_n^{-1/2}\longrightarrow0.
\]
Multiplication by the fixed \(\omega_{\mathrm{win}}\) preserves these limits. Choose
\(n_{\omega_{\mathrm{win}}}\in\mathbb N\) so that, for every \(n\ge n_{\omega_{\mathrm{win}}}\),
\[
\omega_{\mathrm{win}}a_n\le\frac14,\qquad \omega_{\mathrm{win}}b_n\le\frac14,
\]
and define
\[
n_K:=\max\{1,n_{\omega_{\mathrm{win}}},n_E\}.
\]
For \(n\ge n_K\), positivity of \(n,L_n,\omega_{\mathrm{win}}\) ensures
\[
0<\frac{a_n}{\omega_{\mathrm{win}}},\qquad
\omega_{\mathrm{win}}a_n\le\frac14,\qquad
0<\frac{b_n}{\omega_{\mathrm{win}}},\qquad
\omega_{\mathrm{win}}b_n\le\frac14.
\]

At the first threshold, the exact identity from the preceding argument gives
\[
D_n\le C_{\mathrm a}F_n(a_n),
\qquad
F_n(a_n)\le C_{\mathrm a}D_n.
\]
For \(\sigma\in[a_n/\omega_{\mathrm{win}},\omega_{\mathrm{win}}a_n]\), the window bounds therefore imply
\[
\begin{aligned}
D_n
&\le C_{\mathrm a}F_n(a_n)
\le C_{\mathrm a}M_{\mathrm F,\omega_{\mathrm{win}}}F_n(\sigma)
\le C_KF_n(\sigma),\\
F_n(\sigma)
&\le M_{\mathrm F,\omega_{\mathrm{win}}}F_n(a_n)
\le M_{\mathrm F,\omega_{\mathrm{win}}}C_{\mathrm a}D_n
\le C_KD_n.
\end{aligned}
\]
Since \(C_K>0\), these are precisely
\[
C_K^{-1}D_n\le F_n(\sigma)\le C_KD_n.
\]

For \(\sigma\in[b_n/\omega_{\mathrm{win}},\omega_{\mathrm{win}}b_n]\), combine both window comparisons with the second-threshold bounds:
\[
\begin{aligned}
P_n(\sigma)
&\le M_{\mathrm P,\omega_{\mathrm{win}}}P_n(b_n)
\le M_{\mathrm P,\omega_{\mathrm{win}}}C_EF_n(b_n)\\
&\le M_{\mathrm P,\omega_{\mathrm{win}}}C_EM_{\mathrm F,\omega_{\mathrm{win}}}F_n(\sigma)
\le C_KF_n(\sigma),\\
F_n(\sigma)
&\le M_{\mathrm F,\omega_{\mathrm{win}}}F_n(b_n)
\le M_{\mathrm F,\omega_{\mathrm{win}}}C_EP_n(b_n)\\
&\le M_{\mathrm F,\omega_{\mathrm{win}}}C_EM_{\mathrm P,\omega_{\mathrm{win}}}P_n(\sigma)
\le C_KP_n(\sigma).
\end{aligned}
\]
Dividing the first inequality by \(C_K\) establishes
\[
C_K^{-1}P_n(\sigma)\le F_n(\sigma)\le C_KP_n(\sigma).
\]
The choices of \(C_K,n_K\) precede the choice of \(\sigma\), and the same choices apply to both entire windows.
% lean: CausalSmith.Stat.NoisydoseWeakdesignTransition.frontier_transition_windows

\item \textbf{Fixed positive noise and the frontier rate.}
Fix \(\sigma\in(0,1/4]\). Since both thresholds tend to zero, eventually
\[
a_n<\sigma,\qquad b_n<\sigma.
\]
The two branch tests in \cref{obj:def:frontier-rate} then select
\[
b_{n,\sigma}=P_n(\sigma)
=\frac{\log(e+\sigma^2L_n)}{L_n}.
\]

For natural \(n\) with \(\log n\ge1\), define
\[
\zeta_n:=\frac{\log\log n}{\log n}.
\]
Both \(\log n\) and \(\log\log n\) tend to infinity. Hence eventually
\[
\log n\ge1,\qquad
\log\log n\ge0,\qquad
\log\log n\ge-2\log(\sigma^2),
\qquad
\log\log n\ge\log(e+2\sigma^2).
\]
Also \(L_n=1+\log n>0\). Monotonicity of the logarithm gives the lower numerator bound
\[
\begin{aligned}
\log(e+\sigma^2L_n)
&\ge\log(\sigma^2\log n)\\
&=\log(\sigma^2)+\log\log n
\ge\frac12\log\log n.
\end{aligned}
\]
For the upper bound, \(\log n\ge1\) implies
\[
e+\sigma^2L_n
=e+\sigma^2(1+\log n)
\le(e+2\sigma^2)\log n.
\]
Thus
\[
\log(e+\sigma^2L_n)
\le\log(e+2\sigma^2)+\log\log n
\le2\log\log n.
\]
We also have
\[
0\le\zeta_n\le\log\log n.
\]
Consequently,
\[
\frac14\zeta_nL_n
=\frac14\log\log n+\frac14\zeta_n
\le\frac12\log\log n
\le\log(e+\sigma^2L_n),
\]
and
\[
\log(e+\sigma^2L_n)
\le2\log\log n
\le2\log\log n+2\zeta_n
=2\zeta_nL_n.
\]
Dividing by \(L_n>0\) proves
\[
0\le\zeta_n,\qquad
\frac14\zeta_n\le P_n(\sigma)\le2\zeta_n.
\]
Since \(\beta\ge0\), taking the \(\beta\)-th power preserves these comparisons. Together with the selected branch, this gives
\[
0\le\zeta_n^\beta,\qquad
4^{-\beta}\zeta_n^\beta
\le\rho_{n,\sigma}
\le2^\beta\zeta_n^\beta.
\]
Choose \(n_F\in\mathbb N\) so that all these assertions hold for every \(n\ge n_F\).
% lean: CausalSmith.Stat.NoisydoseWeakdesignTransition.fixedNoise_frontierRate_bounds

\item \textbf{Fixed-noise risk and length constants.}
Define
\[
c_\sigma:=c_L4^{-\beta},
\qquad
C_\sigma:=
\max\left\{\max\{C_R,C_H\}2^\beta,c_\sigma\right\}+1,
\]
\[
n_\sigma:=\max\{n_F,n_L,n_R,n_H\}.
\]
Then \(0<c_\sigma<C_\sigma\). For every public path space and every \(n\ge n_\sigma\), the lower-rate comparison and the common lower bound yield
\[
c_\sigma\zeta_n^\beta
\le c_L\rho_{n,\sigma}\le R_n(\sigma),
\qquad
c_\sigma\zeta_n^\beta
\le c_L\rho_{n,\sigma}\le H_n(\sigma).
\]
The upper-rate comparison and the two upper bounds yield
\[
R_n(\sigma)
\le C_R\rho_{n,\sigma}
\le C_R2^\beta\zeta_n^\beta
\le C_\sigma\zeta_n^\beta,
\]
\[
H_n(\sigma)
\le C_H\rho_{n,\sigma}
\le C_H2^\beta\zeta_n^\beta
\le C_\sigma\zeta_n^\beta.
\]
Substituting the definition of \(\zeta_n\) proves the fixed-positive-noise assertions. Every constant and sample threshold was chosen independently of the public path space, as required.
% lean: CausalSmith.Stat.NoisydoseWeakdesignTransition.uniform_frontier
\end{enumerate}
\end{proof}

\begin{proof}[Proof of \cref{obj:prop:error-free-reduction}]
\begin{enumerate}
\item Apply \cref{obj:thm:uniform-frontier}, using the filtered Jacobi localization estimate \citep[Theorem 2.4, equation (2.14)]{PetrushevXu2005}, the local Lipschitz estimate \citep[Theorem 2.2, equation (2.4)]{https://people.math.sc.edu/pencho/Publications/kpx-06-28-06-web.pdf}, Jacobi orthogonality and the squared-norm formula \citep[Table 18.3.1, Jacobi row and caption]{https://dlmf.nist.gov/18.3}, and the gamma-ratio asymptotic \citep[equation (5.11.12)]{https://dlmf.nist.gov/5.11.E12}. It supplies constants and a threshold satisfying
\[
0<c_{\mathrm F}<C_{\mathrm F},
\qquad n_{\mathrm F}\in\mathbb N,
\]
such that, for every public potential-path space, every \(n\ge n_{\mathrm F}\), and every \(\sigma\in[0,1/4]\),
\[
c_{\mathrm F}\rho_{n,\sigma}
\le R_n(\sigma)\le C_{\mathrm F}\rho_{n,\sigma},
\qquad
c_{\mathrm F}\rho_{n,\sigma}
\le H_n(\sigma)\le C_{\mathrm F}\rho_{n,\sigma},
\]
with the prescribed noncoverage probability \(\alpha\in(0,1/2)\). These constants and the threshold are uniform over the public potential-path space.
% lean: CausalSmith.Stat.NoisydoseWeakdesignTransition.uniform_frontier

\item Use the effective exponent from \cref{obj:def:frontier-rate},
\[
d=2\beta+\kappa+1.
\]
Since \(0\le n^{-1/d}\), the first branch of the localization scale at zero error gives
\[
b_{n,0}=n^{-1/d},
\qquad
\rho_{n,0}
=(n^{-1/d})^\beta
=n^{-\beta/d}.
\]
The last identity follows by multiplying the exponents. Substituting \(\sigma=0\in[0,1/4]\) into the bounds above yields
\[
c_{\mathrm F}n^{-\beta/d}
\le R_n(0)\le C_{\mathrm F}n^{-\beta/d},
\qquad
c_{\mathrm F}n^{-\beta/d}
\le H_n(0)\le C_{\mathrm F}n^{-\beta/d}
\]
for every public potential-path space and every \(n\ge n_{\mathrm F}\). This proves the first assertion.
% lean: CausalSmith.Stat.NoisydoseWeakdesignTransition.frontierRate_zero

\item Fix the benchmark parameters, evaluation point, neighborhood, and design distribution in the statement. Define its minimax absolute risk by
\[
\mathfrak R_{\mathrm G,n}
=
\inf_{\widehat f}\sup_f
E_f\bigl|\widehat f(x_*)-f(x_*)\bigr|,
\]
where the functions and measurable estimators range over the benchmark class and procedures specified there. The exponent restrictions imply
\[
0<\beta\le1,
\qquad 0\le\kappa\le2.
\]
Thus Gaiffas's direct Gaussian-response benchmark applies with smoothness \(\beta\), design exponent \(\kappa\), and loss parameter one \citep[Definition 2 and Theorem 1, equations (2.1), (2.4), and (2.5), pp.~3--4]{Gaiffas2005}. It supplies constants and a threshold satisfying
\[
0<c_{\mathrm B}<C_{\mathrm B},
\qquad n_{\mathrm B}\in\mathbb N,
\]
and, for every \(n\ge n_{\mathrm B}\),
\[
c_{\mathrm B}n^{-\beta/d}
\le \mathfrak R_{\mathrm G,n}
\le C_{\mathrm B}n^{-\beta/d}.
\]
This invocation uses the stated local polynomial-approximation class, whose polynomial degree is at most \(\lfloor\beta\rfloor\), including degree at most one when \(\beta=1\). The benchmark constants and threshold depend on the fixed benchmark data.
% lean: CausalSmith.Stat.NoisydoseWeakdesignTransition.GaiffasDirectBenchmark

\item Define the comparison constants and threshold by
\[
c_{\mathrm{cmp}}=\frac{c_{\mathrm F}}{C_{\mathrm B}},
\qquad
C_{\mathrm{cmp}}
=\frac{C_{\mathrm F}}{c_{\mathrm B}}+c_{\mathrm{cmp}}+1,
\qquad
n_{\mathrm{cmp}}=\max\{n_{\mathrm F},n_{\mathrm B}\}.
\]
Positivity of the preceding constants gives
\[
c_{\mathrm{cmp}}>0,
\qquad
C_{\mathrm{cmp}}-c_{\mathrm{cmp}}
=\frac{C_{\mathrm F}}{c_{\mathrm B}}+1>0.
\]
Moreover,
\[
c_{\mathrm{cmp}}C_{\mathrm B}=c_{\mathrm F},
\qquad
C_{\mathrm{cmp}}c_{\mathrm B}
=C_{\mathrm F}+(c_{\mathrm{cmp}}+1)c_{\mathrm B}
\ge C_{\mathrm F}.
\]
For every \(n\ge n_{\mathrm{cmp}}\), both preceding risk bounds apply. Multiplying the benchmark upper bound by \(c_{\mathrm{cmp}}\) gives
\[
c_{\mathrm{cmp}}\mathfrak R_{\mathrm G,n}
\le c_{\mathrm{cmp}}C_{\mathrm B}n^{-\beta/d}
=c_{\mathrm F}n^{-\beta/d}
\le R_n(0).
\]
Conversely, since \(n^{-\beta/d}\ge0\),
\[
R_n(0)
\le C_{\mathrm F}n^{-\beta/d}
\le C_{\mathrm{cmp}}c_{\mathrm B}n^{-\beta/d}
\le C_{\mathrm{cmp}}\mathfrak R_{\mathrm G,n}.
\]
The comparison constants and threshold are uniform over the public potential-path space: the frontier constants have that uniformity, and the benchmark constants depend on the fixed benchmark data. This proves the second assertion.
% lean: CausalSmith.Stat.NoisydoseWeakdesignTransition.error_free_reduction
\end{enumerate}
\end{proof}

\begin{proof}[Proof of \cref{obj:thm:observable-upper}]
\begin{enumerate}
\item Fix public class constants \(K\), \(\beta\in(0,1]\), and \(\kappa\in[0,2]\). The Fourier and polynomial comparison-score conclusions of \cref{obj:lem:dictionary-score-rate} supply positive constants \(C_{\mathrm F},C_{\mathrm M}\) and integer thresholds \(n_{\mathrm F},n_{\mathrm M}\), depending on \(K,\beta,\kappa\), such that
\[
\begin{aligned}
n\ge n_{\mathrm F},\quad \sigma\le L_n^{-1/2}
&\ \Longrightarrow\
\exists k\ge0:\ 
\text{the Fourier pair with index }k\text{ belongs to }\mathcal D_{n,\sigma},
\quad A_q(K)\le C_{\mathrm F}\rho_{n,\sigma},\\
n\ge n_{\mathrm M},\quad L_n^{-1/2}<\sigma
&\ \Longrightarrow\
\exists j\ge0:\ 
\text{the polynomial pair with index }2^j+1\text{ belongs to }\mathcal D_{n,\sigma},
\quad A_q(K)\le C_{\mathrm M}\rho_{n,\sigma},
\end{aligned}
\]
for every \(\sigma\in[0,1/4]\). These inequalities concern public scores and rates determined by \(\beta,\kappa,n,\sigma\), so the constants and thresholds can be fixed uniformly over the potential-path spaces. Set
\[
C=\max\{C_{\mathrm F},C_{\mathrm M}\},
\qquad
n_0=\max\{3,n_{\mathrm F},n_{\mathrm M}\}.
\]
Then \(C>0\). % lean: CausalSmith.Stat.NoisydoseWeakdesignTransition.fourier_dictionary_witness; CausalSmith.Stat.NoisydoseWeakdesignTransition.inverseheat_dictionary_witness

\item Fix a potential-path space, an integer \(n\ge n_0\), and \(\sigma\in[0,1/4]\). In particular,
\[
n\ge3,\qquad n\ge n_{\mathrm F},\qquad n\ge n_{\mathrm M},
\qquad L_n=\log(en)\ge0.
\]
The three expressions defining \(b_{n,\sigma}\) in \cref{obj:def:frontier-rate} are nonnegative: the direct expression is a power of a positive sample size, the Fourier expression has nonnegative numerator and positive denominator, and the polynomial expression has numerator \(\log(e+\sigma^2L_n)\ge0\) and denominator \(L_n>0\). Consequently,
\[
b_{n,\sigma}\ge0,
\qquad
\rho_{n,\sigma}=b_{n,\sigma}^{\beta}\ge0.
\]
This permits enlargement of either comparison constant to \(C\). % lean: CausalSmith.Stat.NoisydoseWeakdesignTransition.observable_upper

\item Suppose \(\sigma\le L_n^{-1/2}\). Choose the Fourier comparison pair from the first witness bound. Its dictionary membership and \(n>0\) allow application of the admissibility conclusion of \cref{obj:lem:dictionary-score-rate}. Thus its functions are measurable and satisfy
\[
\begin{gathered}
q(t)\ge0\quad(t\in[-1/2,1/2]),
\qquad 0<I_\kappa(q)<\infty,\\
E_Z|\ell_q(t+\sigma Z)|<\infty,
\qquad E_Z[\ell_q(t+\sigma Z)]=q(t)
\quad(t\in[-1/2,1/2]),
\qquad V_q<\infty.
\end{gathered}
\]
Moreover,
\[
A_q(K)\le C_{\mathrm F}\rho_{n,\sigma}
\le C\rho_{n,\sigma},
\]
where the last inequality follows from \(C_{\mathrm F}\le C\) and \(\rho_{n,\sigma}\ge0\). This proves the Fourier comparison guarantee, including equality at the cutoff. % lean: CausalSmith.Stat.NoisydoseWeakdesignTransition.observable_upper.hFourier

\item Suppose \(L_n^{-1/2}<\sigma\). Choose the polynomial comparison pair from the second witness bound, with
\[
m=2^j+1,
\qquad
(q,\ell_q)=(q_m^{\mathrm M},\ell_m^{\mathrm M}).
\]
Its dictionary membership again gives, by \cref{obj:lem:dictionary-score-rate}, measurability and
\[
\begin{gathered}
q(t)\ge0\quad(t\in[-1/2,1/2]),
\qquad 0<I_\kappa(q)<\infty,\\
E_Z|\ell_q(t+\sigma Z)|<\infty,
\qquad E_Z[\ell_q(t+\sigma Z)]=q(t)
\quad(t\in[-1/2,1/2]),
\qquad V_q<\infty.
\end{gathered}
\]
Its score satisfies
\[
A_q(K)\le C_{\mathrm M}\rho_{n,\sigma}
\le C\rho_{n,\sigma}.
\]
This proves the polynomial comparison guarantee. % lean: CausalSmith.Stat.NoisydoseWeakdesignTransition.observable_upper.hPoly

\item Let \(\widehat q\) be the selected latent weight in \cref{obj:def:total-estimator}. The fixed-order minimization gives
\[
(\widehat q,\ell_{\widehat q})\in\mathcal D_{n,\sigma},
\qquad
A_{\widehat q}(K)\le A_q(K)
\quad\text{for every }(q,\ell_q)\in\mathcal D_{n,\sigma}.
\]
If \(\sigma\le L_n^{-1/2}\), compare with the Fourier witness; otherwise \(L_n^{-1/2}<\sigma\), and compare with the polynomial witness. In either case,
\[
A_{\widehat q}(K)\le A_q(K)\le C\rho_{n,\sigma}.
\]
% lean: CausalSmith.Stat.NoisydoseWeakdesignTransition.selectedTag_minimizes; CausalSmith.Stat.NoisydoseWeakdesignTransition.observable_upper.hscore

\item Fix \(P\in\mathcal P_{K,\beta,\kappa,\sigma}\) satisfying \cref{obj:ass:iid}. Applied to the selected dictionary pair, \cref{obj:lem:dictionary-score-rate} supplies all its weight-admissibility conditions, finite extended second moment, and the law-specific integrability
\[
E_P\!\left[|\ell_{\widehat q}(W-a_0)|\right]<\infty.
\]
Hence \cref{obj:lem:public-error-certificate} yields
\[
E_{Q_P^n}\!\left[
|\widehat\theta_{\widehat q}-\theta(P)|
\right]\le A_{\widehat q}(K).
\]
For each \(r\in\{0,1,2\}\), the sample index \(i=r\) exists because \(n\ge3\), and \(i\bmod3=r\). Therefore
\[
r\in\mathcal I_r,\qquad
\mathcal I_r\ne\varnothing\quad(r=0,1,2),
\qquad
\widehat\theta=\widehat\theta_{\widehat q}.
\]
Using the total-estimator formula in \cref{obj:def:total-estimator} and the selected-score bound,
\[
E_{Q_P^n}\!\left[|\widehat\theta-\theta(P)|\right]
=
E_{Q_P^n}\!\left[
|\widehat\theta_{\widehat q}-\theta(P)|
\right]
\le A_{\widehat q}(K)
\le C\rho_{n,\sigma}.
\]
The choices of \(C,n_0\) apply to every permitted space, law, and public error scale, as required. % lean: CausalSmith.Stat.NoisydoseWeakdesignTransition.public_error_certificate; CausalSmith.Stat.NoisydoseWeakdesignTransition.blocksNonempty_of_three; CausalSmith.Stat.NoisydoseWeakdesignTransition.observable_upper
\end{enumerate}
\end{proof}

\begin{proof}[Proof of \cref{obj:thm:finite-honesty}]
Fix the parameters, a public potential-outcome space, an allowed error scale \(\sigma\), and a structural law \(P\) satisfying the stated conditions.

\begin{enumerate}
\item The observed-record map is measurable, so its image law is a probability measure. Its \(n\)-fold product \(Q_P^n\), defined in \cref{obj:def:observed-experiment}, is therefore a probability measure, including when \(n=0\).

The finite partition by the two observed strata and \cref{obj:def:conditional-potential-mean} give
\[
\theta(P)
=\sum_{x\in\mathcal X}E_P[\mathbf1\{X=x\}Y(a_0)]
=\sum_{x\in\mathcal X}p_x\mu_x(a_0).
\]
The expectations exist by \cref{obj:ass:bounded-potential-outcomes}, and the conditional means are defined because \cref{obj:ass:stratum-overlap} supplies positive stratum probabilities. Since
\[
p_x\ge0,\qquad \sum_{x\in\mathcal X}p_x=1,
\]
the bounds in \cref{obj:ass:mean-range} yield
\[
0
=0\sum_{x\in\mathcal X}p_x
\le\theta(P)
\le1\sum_{x\in\mathcal X}p_x
=1.
\]
% lean: CausalSmith.Stat.NoisydoseWeakdesignTransition.experiment_isProbabilityMeasure
% lean: CausalSmith.Stat.NoisydoseWeakdesignTransition.causalTarget_mem_Icc

\item First suppose that all three sample blocks in \cref{obj:def:total-estimator} are nonempty. An index in \(\mathcal I_2\) has remainder two modulo three, so
\[
i\in\mathcal I_2
\quad\Longrightarrow\quad
2\le i<n.
\]
Consequently \(n\ge3\).

The selection rule in \cref{obj:def:total-estimator} gives
\[
\widehat q\in\mathcal D_{n,\sigma},
\qquad
A_{\widehat q}(K)\le A_q(K)
\quad\text{for every }q\in\mathcal D_{n,\sigma}.
\]
By the admissibility assertion of \cref{obj:lem:dictionary-score-rate}, the selected pair satisfies the measurability and pointwise inversion conditions, together with
\[
\widehat q(t)\ge0\quad(t\in[-1/2,1/2]),
\qquad
0<I_\kappa(\widehat q)<\infty,
\qquad
V_{\widehat q}^{(16)}<\infty,
\qquad
E_P|\ell_{\widehat q}(V)|<\infty.
\]
Its public denominator and bias bounds have the formulas
\[
b_{\widehat q}(K)
=p_{\min}c_{\mathrm{lo}}I_\kappa(\widehat q)>0,
\qquad
B_{\widehat q}(K)
=\frac{c_{\mathrm{hi}}}{c_{\mathrm{lo}}I_\kappa(\widehat q)}
  \int_{-1/2}^{1/2}|t|^{\kappa+\beta}\widehat q(t)\,dt
\ge0.
\]
Here the bias bound is nonnegative because its integrand is nonnegative. Thus the score formula in \cref{obj:synth_4} gives
\[
A_{\widehat q}(K)
=
B_{\widehat q}(K)
+\frac{12}{b_{\widehat q}(K)}\sqrt{\frac{V_{\widehat q}(K)}{n}}
+2\sqrt{\frac3n}
\ge2\sqrt{\frac3n}>0.
\]
The product experiment satisfies the sampling condition in \cref{obj:ass:iid}. Applying \cref{obj:lem:public-error-certificate} to the selected pair, and using
\(\widehat\theta=\widehat\theta_{\widehat q}\) on nonempty blocks, therefore gives
\[
E_{Q_P^n}\!\left[|\widehat\theta-\theta(P)|\right]
\le A_{\widehat q}(K).
\]
% lean: CausalSmith.Stat.NoisydoseWeakdesignTransition.three_le_of_blocksNonempty
% lean: CausalSmith.Stat.NoisydoseWeakdesignTransition.selectedTag_minimizes
% lean: CausalSmith.Stat.NoisydoseWeakdesignTransition.Aq_pos
% lean: CausalSmith.Stat.NoisydoseWeakdesignTransition.public_error_certificate

\item The stratum estimates are clipped to \([0,1]\). The empirical stratum frequencies are nonnegative and sum to one, since the two strata partition the nonempty frequency block. Hence
\[
\widehat p_x\ge0,\qquad
\sum_{x\in\mathcal X}\widehat p_x=1,\qquad
0\le\widehat\mu_x(\widehat q)\le1,
\]
and
\[
0
\le
\widehat\theta
=\sum_{x\in\mathcal X}\widehat p_x\widehat\mu_x(\widehat q)
\le1.
\]
Together with the target range, this implies
\[
0\le|\widehat\theta-\theta(P)|\le1.
\]
The error is measurable and integrable under the probability law \(Q_P^n\).

Define the measurable event
\[
\mathcal G_{\mathrm{cover}}
=
\left\{\text{data}:
|\widehat\theta-\theta(P)|\le A_{\widehat q}(K)/\alpha\right\}.
\]
Markov's inequality and the certificate give
\[
A_{\widehat q}(K)/\alpha\,
Q_P^n\!\left\{
A_{\widehat q}(K)/\alpha\le|\widehat\theta-\theta(P)|
\right\}
\le
E_{Q_P^n}|\widehat\theta-\theta(P)|
\le A_{\widehat q}(K).
\]
Because \(A_{\widehat q}(K)>0\) and \(\alpha>0\), and the complement of
\(\mathcal G_{\mathrm{cover}}\) is contained in the displayed upper-tail event,
\[
Q_P^n(\mathcal G_{\mathrm{cover}}^c)
\le
Q_P^n\!\left\{
A_{\widehat q}(K)/\alpha\le|\widehat\theta-\theta(P)|
\right\}
\le\alpha,
\qquad
Q_P^n(\mathcal G_{\mathrm{cover}})\ge1-\alpha.
\]
On \(\mathcal G_{\mathrm{cover}}\),
\[
\widehat\theta-A_{\widehat q}(K)/\alpha
\le\theta(P)
\le\widehat\theta+A_{\widehat q}(K)/\alpha.
\]
Combining these inequalities with the target range gives
\[
\max\!\left\{0,\widehat\theta-A_{\widehat q}(K)/\alpha\right\}
\le\theta(P)
\le
\min\!\left\{1,\widehat\theta+A_{\widehat q}(K)/\alpha\right\}.
\]
Thus \cref{obj:def:honest-interval} implies
\[
Q_P^n\{\theta(P)\in C_{n,\sigma}\}
\ge Q_P^n(\mathcal G_{\mathrm{cover}})
\ge1-\alpha.
\]
% lean: CausalSmith.Stat.NoisydoseWeakdesignTransition.totalEstimator_mem_Icc
% lean: mul_meas_ge_le_integral_of_nonneg
% lean: CausalSmith.Stat.NoisydoseWeakdesignTransition.mem_honestInterval_of_abs_le

\item If some sample block is empty, \cref{obj:def:honest-interval} returns the full target range. Therefore
\[
C_{n,\sigma}=[0,1],
\qquad
Q_P^n\{\theta(P)\in C_{n,\sigma}\}=1\ge1-\alpha.
\]
This case includes \(n=0\), and completes the coverage assertion for every integer \(n\ge0\).
% lean: CausalSmith.Stat.NoisydoseWeakdesignTransition.finite_honesty

\item For the length assertion, choose a constant \(C_{\mathrm{up}}>0\) and a threshold \(n_{\mathrm{up}}\in\mathbb N\) supplied by \cref{obj:thm:observable-upper}. Set
\[
C=\frac{2}{\alpha}C_{\mathrm{up}},
\qquad
n_0=\max\{3,n_{\mathrm{up}}\}.
\]
Fix \(n\ge n_0\). For each block index \(r\in\{0,1,2\}\),
\[
r<n,\qquad r\bmod3=r,\qquad r\in\mathcal I_r,
\]
so all blocks are nonempty.

Write
\[
L_n=\log(en),
\]
as in \cref{obj:def:frontier-rate}. If \(\sigma\le L_n^{-1/2}\), the Fourier comparison in \cref{obj:thm:observable-upper} supplies a dictionary pair whose score satisfies
\[
A_q(K)\le C_{\mathrm{up}}\rho_{n,\sigma},
\qquad q\in\mathcal D_{n,\sigma}.
\]
The minimizing property consequently gives
\[
A_{\widehat q}(K)\le A_q(K)\le C_{\mathrm{up}}\rho_{n,\sigma}.
\]
If \(L_n^{-1/2}<\sigma\), the polynomial comparison supplies a dictionary pair with
\[
A_q(K)\le C_{\mathrm{up}}\rho_{n,\sigma},
\qquad q\in\mathcal D_{n,\sigma},
\]
and again
\[
A_{\widehat q}(K)\le A_q(K)\le C_{\mathrm{up}}\rho_{n,\sigma}.
\]
These two cases exhaust the allowed error scales. Dictionary admissibility and the score formula also give
\[
A_{\widehat q}(K)\ge2\sqrt{\frac3n}>0.
\]
% lean: CausalSmith.Stat.NoisydoseWeakdesignTransition.observable_upper
% lean: CausalSmith.Stat.NoisydoseWeakdesignTransition.blocksNonempty_of_three
% lean: CausalSmith.Stat.NoisydoseWeakdesignTransition.selectedTag_minimizes
% lean: CausalSmith.Stat.NoisydoseWeakdesignTransition.Aq_pos

\item Define the interval endpoints by
\[
c_{\mathrm{left}}
=\max\!\left\{0,\widehat\theta-\frac{A_{\widehat q}(K)}{\alpha}\right\},
\qquad
c_{\mathrm{right}}
=\min\!\left\{1,\widehat\theta+\frac{A_{\widehat q}(K)}{\alpha}\right\}.
\]
They satisfy
\[
c_{\mathrm{left}}\ge\widehat\theta-\frac{A_{\widehat q}(K)}{\alpha},
\qquad
c_{\mathrm{right}}\le\widehat\theta+\frac{A_{\widehat q}(K)}{\alpha},
\qquad
c_{\mathrm{right}}-c_{\mathrm{left}}\le\frac{2A_{\widehat q}(K)}{\alpha}.
\]
The length formula for a closed interval, together with
\(A_{\widehat q}(K)\ge0\), therefore yields the pointwise bound
\[
\operatorname{length}(C_{n,\sigma})
=\max\{c_{\mathrm{right}}-c_{\mathrm{left}},0\}
\le\frac{2A_{\widehat q}(K)}{\alpha}.
\]
Since the selected score is deterministic and \(Q_P^n\) is a probability measure,
\[
\begin{aligned}
E_{Q_P^n}\!\left[\operatorname{length}(C_{n,\sigma})\right]
&\le E_{Q_P^n}\!\left[\frac{2A_{\widehat q}(K)}{\alpha}\right]\\
&=\frac{2A_{\widehat q}(K)}{\alpha}\\
&\le\frac{2C_{\mathrm{up}}}{\alpha}\rho_{n,\sigma}
=C\rho_{n,\sigma}.
\end{aligned}
\]
The chosen constant and threshold depend only on \(K,\beta,\kappa\), and \(\alpha\), and the argument applies simultaneously to every permitted schedule space, error scale, and model law.
% lean: CausalSmith.Stat.NoisydoseWeakdesignTransition.honestInterval_volume_le
% lean: CausalSmith.Stat.NoisydoseWeakdesignTransition.finite_honesty
\end{enumerate}
\end{proof}

\begin{proof}[Proof of \cref{obj:thm:weighted-packet-construction}]
Fix \(\kappa\in[0,2]\). We use the filter, taper parameter, kernels, and packet from \cref{obj:def:packet-handle}. All constants below are independent of the integer \(m\ge4\).

\begin{enumerate}
\item \emph{Filter properties and cited kernel estimates.}
The standard smooth function
\[
\vartheta(v)=
\begin{cases}
e^{-1/v},&v>0,\\
0,&v\le0
\end{cases}
\qquad(v\in\mathbb R)
\]
gives
\[
\eta(s)=\vartheta\bigl((s-9/8)(15/8-s)\bigr).
\]
Consequently, \(\eta\) is smooth, \(0\le\eta\le1\), and
\[
\eta(s)\ne0\ \Longleftrightarrow\ 9/8<s<15/8.
\]
In particular,
\[
\eta(j/m)\ne0\quad\Longrightarrow\quad m<j<2m.
\]
On the central spectral interval,
\[
(s-9/8)(15/8-s)\ge\frac1{64}
\quad(5/4\le s\le7/4),
\]
so
\[
\eta(s)\ge c_\eta,\qquad c_\eta:=e^{-64}>0.
\]

For Jacobi parameters \(a_{\mathrm J},b_{\mathrm J}>-1/2\), define
\[
\mathcal K_m^{a_{\mathrm J},b_{\mathrm J}}(z_1,z_2)
:=
\sum_{j=0}^{2m}\eta(j/m)
\frac{\mathsf J_j^{(a_{\mathrm J},b_{\mathrm J})}(z_1)
      \mathsf J_j^{(a_{\mathrm J},b_{\mathrm J})}(z_2)}
     {H_j^{a_{\mathrm J},b_{\mathrm J}}}
\qquad(z_1,z_2\in\mathbb R),
\]
and, for \(z,z_1,z_2\in[-1,1]\), define
\[
\mathcal W_{a_{\mathrm J},b_{\mathrm J}}(m;z)
:=
(1-z+m^{-2})^{a_{\mathrm J}+1/2}
(1+z+m^{-2})^{b_{\mathrm J}+1/2},
\qquad
\delta_{\mathrm{ang}}(z_1,z_2)
:=
|\arccos z_1-\arccos z_2|.
\]
The spectral support permits truncation at degree \(2m\). Petrushev--Xu use squared norms normalized by the total Jacobi weight mass
\citep[pp.~1--2, equation~(1.2)]{PetrushevXu2005}. In the present notation, their norms and filtered kernel are
\[
\mathsf h_{\mathrm{PX},j}^{a_{\mathrm J},b_{\mathrm J}}
:=\frac{H_j^{a_{\mathrm J},b_{\mathrm J}}}
        {H_0^{a_{\mathrm J},b_{\mathrm J}}}
\qquad(j\ge0),
\]
\[
\begin{split}
\mathcal K_{\mathrm{PX},m}^{a_{\mathrm J},b_{\mathrm J}}(z_1,z_2)
&:=\sum_{j=0}^{2m}\eta(j/m)
\frac{\mathsf J_j^{(a_{\mathrm J},b_{\mathrm J})}(z_1)
      \mathsf J_j^{(a_{\mathrm J},b_{\mathrm J})}(z_2)}
     {\mathsf h_{\mathrm{PX},j}^{a_{\mathrm J},b_{\mathrm J}}}\\
&=H_0^{a_{\mathrm J},b_{\mathrm J}}
  \mathcal K_m^{a_{\mathrm J},b_{\mathrm J}}(z_1,z_2)
\qquad(z_1,z_2\in\mathbb R).
\end{split}
\]
Since \(H_0^{a_{\mathrm J},b_{\mathrm J}}>0\) depends only on the fixed Jacobi parameters, dividing their localization estimate by this factor absorbs its reciprocal into the localization constant below. The localization estimate of
\citep[Theorem~2.4, equation~(2.14)]{PetrushevXu2005}
provides, for every \(\tau>0\), a positive constant
\(C_{\mathrm{loc},\tau}^{a_{\mathrm J},b_{\mathrm J}}\) such that
\[
|\mathcal K_m^{a_{\mathrm J},b_{\mathrm J}}(z_1,z_2)|
\le
\frac{C_{\mathrm{loc},\tau}^{a_{\mathrm J},b_{\mathrm J}}m}
{\sqrt{\mathcal W_{a_{\mathrm J},b_{\mathrm J}}(m;z_1)
       \mathcal W_{a_{\mathrm J},b_{\mathrm J}}(m;z_2)}
 (1+m\delta_{\mathrm{ang}}(z_1,z_2))^\tau}.
\]
The local angular estimate of
\citep[Theorem~2.2, equation~(2.4)]{https://people.math.sc.edu/pencho/Publications/kpx-06-28-06-web.pdf},
with its fixed neighborhood constant equal to \(1\), provides a positive
\(C_{\mathrm{lip},\tau}^{a_{\mathrm J},b_{\mathrm J}}\) such that
\[
\begin{split}
&|\mathcal K_m^{a_{\mathrm J},b_{\mathrm J}}(z_1,z_2)
 -\mathcal K_m^{a_{\mathrm J},b_{\mathrm J}}(\xi,z_2)|\\
&\qquad\le
\frac{C_{\mathrm{lip},\tau}^{a_{\mathrm J},b_{\mathrm J}}m^2
      \delta_{\mathrm{ang}}(z_1,\xi)}
{\sqrt{\mathcal W_{a_{\mathrm J},b_{\mathrm J}}(m;z_2)
       \mathcal W_{a_{\mathrm J},b_{\mathrm J}}(m;\zeta)}
 (1+m\delta_{\mathrm{ang}}(z_2,\zeta))^\tau},
\end{split}
\]
whenever \(z_1,z_2,\xi,\zeta\in[-1,1]\) and
\[
\delta_{\mathrm{ang}}(z_1,\xi)\le m^{-1},
\qquad
\delta_{\mathrm{ang}}(\zeta,\xi)\le m^{-1}.
\]
The parameter and filter hypotheses hold for both Jacobi systems used below.
% lean: CausalSmith.Stat.NoisydoseWeakdesignTransition.FilteredJacobiLipschitz

\item \emph{Norms and uniform gamma-ratio bounds.}
Jacobi orthogonality and the positive squared-norm formula are supplied by
\citep[Table~18.3.1, Jacobi row and caption]{https://dlmf.nist.gov/18.3}:
\[
H_j^{a_{\mathrm J},b_{\mathrm J}}
=
\frac{2^{a_{\mathrm J}+b_{\mathrm J}+1}}
     {2j+a_{\mathrm J}+b_{\mathrm J}+1}
\frac{\Gamma(j+a_{\mathrm J}+1)\Gamma(j+b_{\mathrm J}+1)}
     {\Gamma(j+1)\Gamma(j+a_{\mathrm J}+b_{\mathrm J}+1)}
>0.
\]
Orthogonality applies against every polynomial of degree less than \(j\).

For fixed positive shifts \(\lambda_1,\lambda_2\), the asymptotic in
\citep[Equation~5.11.12]{https://dlmf.nist.gov/5.11.E12}
gives an integer cutoff \(j_\Gamma\) such that
\[
\frac12
\le
\frac{\Gamma(j+\lambda_1)}
     {\Gamma(j+\lambda_2)j^{\lambda_1-\lambda_2}}
\le2
\qquad(j\ge j_\Gamma,\ j\ge1).
\]
All these normalized ratios are positive. Thus the constants
\[
\begin{split}
c_\Gamma(\lambda_1,\lambda_2)
&:=
\min\left(
\{1/2\}\cup
\left\{
\frac{\Gamma(j+\lambda_1)}
     {\Gamma(j+\lambda_2)j^{\lambda_1-\lambda_2}}
:1\le j<j_\Gamma
\right\}\right),\\
C_\Gamma(\lambda_1,\lambda_2)
&:=
\max\left(
\{2\}\cup
\left\{
\frac{\Gamma(j+\lambda_1)}
     {\Gamma(j+\lambda_2)j^{\lambda_1-\lambda_2}}
:1\le j<j_\Gamma
\right\}\right)
\end{split}
\]
are positive and satisfy
\[
c_\Gamma(\lambda_1,\lambda_2)j^{\lambda_1-\lambda_2}
\le
\frac{\Gamma(j+\lambda_1)}{\Gamma(j+\lambda_2)}
\le
C_\Gamma(\lambda_1,\lambda_2)j^{\lambda_1-\lambda_2}
\qquad(j\ge1).
\]
The singleton sets in these definitions also cover an empty initial segment.
% lean: CausalSmith.Stat.NoisydoseWeakdesignTransition.gamma_ratio_integer_power_bounds

\item \emph{Endpoint normalization for \(\kappa>0\).}
Suppose first that \(\kappa>0\), and set
\[
b:=\frac{\kappa-1}{2}>-\frac12.
\]
The standard endpoint evaluation, followed by the gamma recurrence, gives
\[
\mathsf J_j^{(4,b)}(-1)
=
\frac{(-1)^j}{j!}\prod_{i=0}^{j-1}(b+1+i)
=
(-1)^j\frac{\Gamma(j+b+1)}
{\Gamma(b+1)\Gamma(j+1)}
\qquad(j\ge0),
\]
where the empty product is \(1\). Every factor in the product is positive, so the endpoint value is nonzero. Combining this identity with the norm formula yields
\[
\frac{\mathsf J_j^{(4,b)}(-1)^2}{H_j^{4,b}}
=
\frac{2j+b+5}{2^{b+5}\Gamma(b+1)^2}
\frac{\Gamma(j+b+1)}{\Gamma(j+1)}
\frac{\Gamma(j+b+5)}{\Gamma(j+5)}.
\]
For \(j\ge1\),
\[
2j\le2j+b+5\le(b+7)j.
\]
Applying the preceding gamma bounds to the shift pairs
\((b+1,1)\) and \((b+5,5)\), define
\[
\begin{split}
c_{\mathrm{spec,end}}
&:=
\frac{2c_\Gamma(b+1,1)c_\Gamma(b+5,5)}
     {2^{b+5}\Gamma(b+1)^2},\\
C_{\mathrm{spec,end}}
&:=
\frac{(b+7)C_\Gamma(b+1,1)C_\Gamma(b+5,5)}
     {2^{b+5}\Gamma(b+1)^2}.
\end{split}
\]
These constants are positive, and \(2b+1=\kappa\) gives
\[
c_{\mathrm{spec,end}}j^\kappa
\le
\frac{\mathsf J_j^{(4,b)}(-1)^2}{H_j^{4,b}}
\le
C_{\mathrm{spec,end}}j^\kappa
\qquad(j\ge1).
\]

Every summand of \(K_m(-1)\) is nonnegative. Strict positivity follows already from the degree
\[
\nu_{\mathrm{pos}}(m):=m+\lfloor m/2\rfloor,
\qquad
\frac98<\frac{\nu_{\mathrm{pos}}(m)}m<\frac{15}8
\quad(m\ge4),
\]
whose filter value and squared endpoint coefficient are positive.

For the quantitative lower bound, use the distinct degrees
\[
\nu_{\mathrm{end}}(m,i)
:=
m+\lfloor m/4\rfloor+1+i,
\qquad
0\le i<\lfloor m/4\rfloor.
\]
They satisfy
\[
\frac54\le\frac{\nu_{\mathrm{end}}(m,i)}m\le\frac74,
\qquad
m\le\nu_{\mathrm{end}}(m,i)\le2m.
\]
Moreover,
\[
m\le4\lfloor m/4\rfloor+3
\le7\lfloor m/4\rfloor
\qquad(m\ge4).
\]
Each selected summand is therefore at least
\(c_\eta c_{\mathrm{spec,end}}m^\kappa\). For the upper bound, inactive summands vanish, whereas each active summand is at most
\(C_{\mathrm{spec,end}}2^\kappa m^\kappa\); there are at most \(2m+1\le3m\) summands. Hence, with
\[
c_{\mathrm{end}}
:=\frac{c_\eta c_{\mathrm{spec,end}}}{7}>0,
\qquad
C_{\mathrm{end}}
:=3\,2^\kappa C_{\mathrm{spec,end}}>0,
\]
we obtain
\[
0<c_{\mathrm{end}}m^{\kappa+1}
\le K_m(-1)\le C_{\mathrm{end}}m^{\kappa+1}
\qquad(m\ge4).
\]
% lean: CausalSmith.Stat.NoisydoseWeakdesignTransition.packetKernel_endpoint_power_bounds

\item \emph{Endpoint localization.}
Define the polynomial coordinate and taper base on the whole line by
\[
z_{\mathrm{end}}(u):=2u^2-1,
\qquad
\Delta(u):=1-u^2
\qquad(u\in\mathbb R).
\]
For \(|u|\le1\), define
\[
\mathcal B_{\mathrm{end}}(m,u)
:=
\sqrt{\mathcal W_{4,b}(m;-1)
      \mathcal W_{4,b}(m;z_{\mathrm{end}}(u))}.
\]
Here \(0\le\Delta(u)\le1\), and
\[
\begin{split}
\mathcal W_{4,b}(m;-1)
&=(2+m^{-2})^{9/2}m^{-\kappa}
\ge m^{-\kappa},\\
\mathcal W_{4,b}(m;z_{\mathrm{end}}(u))
&=(2\Delta(u)+m^{-2})^{9/2}
  (2u^2+m^{-2})^{\kappa/2}.
\end{split}
\]
The inequalities
\[
\Delta(u)^8
\le\Delta(u)^{9/2}
\le(2\Delta(u)+m^{-2})^{9/2},
\qquad
m^{-\kappa}\le(2u^2+m^{-2})^{\kappa/2}
\]
therefore imply
\[
\mathcal B_{\mathrm{end}}(m,u)
\ge\Delta(u)^4m^{-\kappa}.
\]
Also,
\[
\delta_{\mathrm{ang}}(-1,z_{\mathrm{end}}(u))
=2\arcsin|u|
\ge|u|.
\]
Indeed, the equality follows from
\(\arccos(2u^2-1)=\pi-2\arcsin|u|\), and the inequality follows from
\(\sin v\le v\) for \(v\ge0\).

Use localization with \(\tau=\kappa+2\). On \([-1,1]\), positivity of the normalizer gives
\[
\begin{split}
|\psi_m(u)|
&=
\frac{\Delta(u)^4|K_m(z_{\mathrm{end}}(u))|}{K_m(-1)}\\
&\le
\frac{\Delta(u)^4 C_{\mathrm{loc},\kappa+2}^{4,b}m}
{\mathcal B_{\mathrm{end}}(m,u)
 (1+m\delta_{\mathrm{ang}}(-1,z_{\mathrm{end}}(u)))^{\kappa+2}
 K_m(-1)}\\
&\le
C_{\mathrm{dec,end}}(1+m|u|)^{-(\kappa+2)},
\end{split}
\qquad
C_{\mathrm{dec,end}}
:=
\frac{C_{\mathrm{loc},\kappa+2}^{4,b}}{c_{\mathrm{end}}}>0.
\]
Outside this interval \(\psi_m=0\), so the same bound holds for every real \(u\).
% lean: CausalSmith.Stat.NoisydoseWeakdesignTransition.packetPsi_endpoint_localization

\item \emph{Endpoint derivative and continuous differentiability.}
Fix \(0<|u|<1\). In the local angular estimate take the fixed second argument to be \(-1\), the reference and auxiliary arguments both to be \(z_{\mathrm{end}}(u)\), and the varying argument to be \(z_{\mathrm{end}}(v)\). For \(v\) sufficiently close to \(u\), all arguments belong to \([-1,1]\) and the required angular distance is at most \(m^{-1}\). Symmetry of the finite kernel sum then gives
\[
\begin{split}
&|K_m(z_{\mathrm{end}}(v))-K_m(z_{\mathrm{end}}(u))|\\
&\quad\le
\frac{C_{\mathrm{lip},1}^{4,b}m^2}
{\mathcal B_{\mathrm{end}}(m,u)
 (1+m\delta_{\mathrm{ang}}(-1,z_{\mathrm{end}}(u)))}
|\arccos z_{\mathrm{end}}(v)-\arccos z_{\mathrm{end}}(u)|.
\end{split}
\]
Dividing by \(|v-u|\) and passing to the limit is valid because the kernel is a finite polynomial sum. The identities
\[
1-z_{\mathrm{end}}(u)^2=4u^2\Delta(u),
\qquad
\left|\frac{d}{du}\arccos z_{\mathrm{end}}(u)\right|
=\frac2{\sqrt{\Delta(u)}}
\]
yield
\[
|(K_m\circ z_{\mathrm{end}})'(u)|
\le
\frac{2C_{\mathrm{lip},1}^{4,b}m^2}
{\mathcal B_{\mathrm{end}}(m,u)
 (1+m\delta_{\mathrm{ang}}(-1,z_{\mathrm{end}}(u)))
 \sqrt{\Delta(u)}}
\le
\frac{2C_{\mathrm{lip},1}^{4,b}m^2}
{\mathcal B_{\mathrm{end}}(m,u)\sqrt{\Delta(u)}}.
\]
Localization with exponent \(1\) likewise gives
\[
|K_m(z_{\mathrm{end}}(u))|
\le
\frac{C_{\mathrm{loc},1}^{4,b}m}
{\mathcal B_{\mathrm{end}}(m,u)}.
\]

For the derivative assembly, use the cubed taper. Since
\[
\Delta(u)^6
\le\Delta(u)^{9/2}
\le(2\Delta(u)+m^{-2})^{9/2},
\]
the same weight comparison as before gives
\[
\frac{\Delta(u)^3}{\mathcal B_{\mathrm{end}}(m,u)}
\le m^\kappa.
\]
Furthermore,
\[
\frac{\Delta(u)^4}{\sqrt{\Delta(u)}}\le\Delta(u)^3,
\]
because \(0<\Delta(u)\le1\). The product rule gives
\[
\psi_m'(u)=
\frac{-8u\Delta(u)^3K_m(z_{\mathrm{end}}(u))
+\Delta(u)^4(K_m\circ z_{\mathrm{end}})'(u)}
{K_m(-1)}.
\]
Using \(|u|\le1\), the preceding inequalities, and the normalization lower bound,
\[
\begin{split}
|\psi_m'(u)|
&\le
\frac{
8C_{\mathrm{loc},1}^{4,b}m
 \Delta(u)^3/\mathcal B_{\mathrm{end}}(m,u)
+
2C_{\mathrm{lip},1}^{4,b}m^2
 \Delta(u)^4/(\mathcal B_{\mathrm{end}}(m,u)\sqrt{\Delta(u)})
}{K_m(-1)}\\
&\le
\frac{8C_{\mathrm{loc},1}^{4,b}m^{\kappa+1}
      +2C_{\mathrm{lip},1}^{4,b}m^{\kappa+2}}
     {c_{\mathrm{end}}m^{\kappa+1}}\\
&\le C_{\mathrm{der,end}}m,
\end{split}
\qquad
C_{\mathrm{der,end}}
:=
\frac{8C_{\mathrm{loc},1}^{4,b}
      +2C_{\mathrm{lip},1}^{4,b}}{c_{\mathrm{end}}}>0.
\]

To obtain continuous differentiability on the whole line, define
\[
\mathcal T(u):=
\begin{cases}
(1-u^2)^4,&|u|\le1,\\
0,&|u|>1.
\end{cases}
\]
For every \(u\in\mathbb R\),
\[
\mathcal T(u)
=
\frac{(1-u^2)^4+(1-u^2)|1-u^2|^3}{2}.
\]
The function \(v\mapsto|v|^3\) is continuously differentiable, with derivative \(3v|v|\). Thus \(\mathcal T\) is continuously differentiable, and
\[
\psi_m(u)
=
\mathcal T(u)\frac{K_m(z_{\mathrm{end}}(u))}{K_m(-1)}
\qquad(u\in\mathbb R)
\]
is a product of continuously differentiable functions. Outside \([-1,1]\) the packet is locally zero, so
\[
\psi_m'(u)=0\qquad(|u|>1).
\]
Continuity of \(\psi_m'\) extends the derivative bound from the complement of
\(\{-1,0,1\}\) to all of \(\mathbb R\).
% lean: CausalSmith.Stat.NoisydoseWeakdesignTransition.packetPsi_endpoint_derivative_bound

\item \emph{Endpoint mass bounds and common constant.}
Set
\[
C_{\mathrm{env,end}}
:=
\max\{C_{\mathrm{dec,end}},C_{\mathrm{der,end}}\},
\qquad
C_{\mathrm{mid,end}}:=2C_{\mathrm{env,end}}.
\]
Both constants are positive. Since the coordinate \(z_{\mathrm{end}}\), taper, and support are invariant under reflection,
\[
\psi_m(-u)=\psi_m(u)\qquad(u\in\mathbb R).
\]
The envelope and \(1+m|u|\ge1\) give
\[
|\psi_m(u)|\le C_{\mathrm{env,end}}\le C_{\mathrm{mid,end}},
\qquad
|\psi_m'(u)|\le C_{\mathrm{mid,end}}m.
\]

For \(0\le u\le1\), the inequality
\((mu)^\kappa\le(1+mu)^\kappa\) gives
\[
u^\kappa(1+mu)^{-(\kappa+2)}
\le m^{-\kappa}(1+mu)^{-2}.
\]
The quadratic envelope has the explicit integral
\[
\int_0^1(1+mu)^{-2}\,du
=
\left[-\frac1{m(1+mu)}\right]_0^1
=
\frac1m-\frac1{m(1+m)}
\le\frac1m.
\]
By evenness,
\[
\begin{split}
\int_{-1}^1|u|^\kappa|\psi_m(u)|\,du
&=
2\int_0^1u^\kappa|\psi_m(u)|\,du\\
&\le
2C_{\mathrm{env,end}}m^{-\kappa}
\int_0^1(1+mu)^{-2}\,du\\
&\le C_{\mathrm{mid,end}}m^{-\kappa-1}.
\end{split}
\]
Continuity ensures integrability of these functions on the compact interval. Pointwise,
\[
\psi_m(u)^2
=
|\psi_m(u)|^2
\le C_{\mathrm{mid,end}}|\psi_m(u)|,
\]
and hence
\[
\int_{-1}^1|u|^\kappa\psi_m(u)^2\,du
\le
C_{\mathrm{mid,end}}
\int_{-1}^1|u|^\kappa|\psi_m(u)|\,du
\le C_{\mathrm{mid,end}}^2m^{-\kappa-1}.
\]
Define
\[
C_{\mathrm{fin,end}}
:=
\max\{C_{\mathrm{mid,end}},C_{\mathrm{mid,end}}^2\}>0.
\]
The supremum, derivative, absolute-mass, and squared-mass bounds all hold with this common constant.
% lean: CausalSmith.Stat.NoisydoseWeakdesignTransition.packetendpoint_bounds

\item \emph{Endpoint value, support, and moments.}
The definition and positive normalizer give
\[
\psi_m(u)=0\quad(|u|>1),
\qquad
\psi_m(0)=\frac{K_m(-1)}{K_m(-1)}=1.
\]

Let \(s\) be an integer with \(0\le s\le m\). Every active spectral degree exceeds \(m\), while the polynomial
\(((1+z)/2)^s\) has degree \(s\). Termwise Jacobi orthogonality therefore gives
\[
\int_{-1}^1
(1-z)^4(1+z)^bK_m(z)
\left(\frac{1+z}{2}\right)^s\,dz=0.
\]
For \(0<u<1\), the substitution \(z=2u^2-1\) has Jacobian \(4u\), and
\[
\begin{split}
&(1-(2u^2-1))^4(1+(2u^2-1))^b
\left(\frac{1+(2u^2-1)}2\right)^s(4u)\\
&\qquad=
2^{6+b}u^\kappa u^{2s}(1-u^2)^4.
\end{split}
\]
Consequently,
\[
\begin{split}
\int_0^1u^\kappa u^{2s}\psi_m(u)\,du
&=
\frac1{2^{6+b}K_m(-1)}
\int_{-1}^1
(1-z)^4(1+z)^bK_m(z)
\left(\frac{1+z}{2}\right)^s\,dz\\
&=0.
\end{split}
\]
The integrand of the full even moment is even, so reflection gives
\[
\int_{-1}^1|u|^\kappa u^{2s}\psi_m(u)\,du
=
2\int_0^1u^\kappa u^{2s}\psi_m(u)\,du=0.
\]
For an odd integer \(j\), reflection changes the sign of the integrand and gives
\[
\int_{-1}^1|u|^\kappa u^j\psi_m(u)\,du
=
-\int_{-1}^1|u|^\kappa u^j\psi_m(u)\,du
=0.
\]
If \(0\le j<m\) is even, writing \(j=2s\) gives \(s\le m\), so its moment vanishes as well. Thus the positive-exponent branch satisfies all assertions with
\[
C:=C_{\mathrm{fin,end}},
\qquad c:=1.
\]
% lean: CausalSmith.Stat.NoisydoseWeakdesignTransition.packetendpoint_bounds

\item \emph{Interior spectral lower bound.}
In the remaining case, \(0\le\kappa\) and \(\kappa\le0\), hence \(\kappa=0\). We use the symmetric Jacobi system with parameters \((4,4)\).

First, the norm formula and gamma bounds for the shift pairs \((5,1)\) and \((5,9)\) give
\[
\begin{split}
H_j^{4,4}
&=
\frac{512}{2j+9}
\frac{\Gamma(j+5)}{\Gamma(j+1)}
\frac{\Gamma(j+5)}{\Gamma(j+9)}\\
&\le
\frac{512}{2j+9}
C_\Gamma(5,1)j^4 C_\Gamma(5,9)j^{-4}
\le\frac{C_{\mathrm{norm,int}}}{j}
\qquad(j\ge1),
\end{split}
\]
where
\[
C_{\mathrm{norm,int}}
:=512C_\Gamma(5,1)C_\Gamma(5,9)>0.
\]

The binomial representation in \cref{obj:lem:jacobi-binomial}, evaluated at zero, gives
\[
\mathsf J_{2s}^{(4,4)}(0)
=
4^{-s}
\sum_{i=0}^{2s}
\binom{2s+4}{i}\binom{2s+4}{2s-i}(-1)^{2s-i}.
\]
The sum is the coefficient of degree \(2s\) in
\((1+z)^{2s+4}(1-z)^{2s+4}=(1-z^2)^{2s+4}\). Therefore
\[
\mathsf J_{2s}^{(4,4)}(0)
=
(-1)^s4^{-s}\binom{2s+4}{s}
\qquad(s\ge0).
\]
In particular, these even-degree center values are nonzero.

Wallis' inequality, in its finite-product form, states
\[
\prod_{i=1}^s\frac{(2i)^2}{(2i-1)(2i+1)}
=
\frac{2^{4s}(s!)^4}{((2s)!)^2(2s+1)}
\le\frac\pi2
\qquad(s\ge0).
\]
The factorial identity
\[
\frac{2^{4s}(s!)^4}{((2s)!)^2(2s+1)}
\left(4^{-s}\binom{2s}{s}\right)^2(2s+1)=1
\]
then gives, for \(s\ge1\),
\[
\left(4^{-s}\binom{2s}{s}\right)^2
\ge\frac2{\pi(2s+1)}
\ge\frac2{3\pi s}.
\]
Since \(\binom{2s+4}{s}\ge\binom{2s}{s}\),
\[
\mathsf J_{2s}^{(4,4)}(0)^2\ge\frac2{3\pi s}.
\]
Together with
\[
0<H_{2s}^{4,4}
\le\frac{C_{\mathrm{norm,int}}}{2s}
\le\frac{C_{\mathrm{norm,int}}}{s},
\]
this proves
\[
\frac{\mathsf J_{2s}^{(4,4)}(0)^2}{H_{2s}^{4,4}}
\ge c_{\mathrm{spec,int}},
\qquad
c_{\mathrm{spec,int}}
:=\frac2{3\pi C_{\mathrm{norm,int}}}>0
\qquad(s\ge1).
\]
% lean: CausalSmith.Stat.NoisydoseWeakdesignTransition.jacobi_four_even_spectral_lower

\item \emph{Interior normalization, including the smallest degrees.}
Define
\[
N_{\mathrm{even}}(m):=\max\{1,\lfloor m/8\rfloor\},
\qquad
\nu_{\mathrm{int}}(m,i)
:=2\bigl(\lceil5m/8\rceil+i\bigr),
\qquad
0\le i<N_{\mathrm{even}}(m).
\]
These are distinct positive even degrees satisfying
\[
\frac54\le\frac{\nu_{\mathrm{int}}(m,i)}m\le\frac74,
\qquad
\nu_{\mathrm{int}}(m,i)\le2m.
\]
For \(m\ge8\), the upper interval bound follows from
\[
\lceil5m/8\rceil+i
\le\frac{5m+7}{8}+\lfloor m/8\rfloor-1
\le\frac{6m-1}{8}\le\frac{7m}{8}.
\]
For \(m=4,5,6,7\), the construction selects respectively the single degree
\(6,8,8,10\), each in the stated interval.

The count satisfies
\[
N_{\mathrm{even}}(m)\ge\frac m{15}\qquad(m\ge4).
\]
Indeed, for \(m\ge8\),
\[
m\le8\lfloor m/8\rfloor+7
\le15\lfloor m/8\rfloor,
\]
and for \(4\le m\le7\) the count is \(1\).

Every summand of \(\widetilde K_m(0)\) is nonnegative. Each selected even degree contributes at least
\(c_\eta c_{\mathrm{spec,int}}\). Hence
\[
\widetilde K_m(0)
\ge
N_{\mathrm{even}}(m)c_\eta c_{\mathrm{spec,int}}
\ge c_{\mathrm{int}}m>0,
\qquad
c_{\mathrm{int}}
:=\frac{c_\eta c_{\mathrm{spec,int}}}{15}>0.
\]
% lean: CausalSmith.Stat.NoisydoseWeakdesignTransition.packetKernel_interior_lower_of_spectral_lower

\item \emph{Interior localization and derivative.}
For \(|u|\le1\), define
\[
\mathcal B_{\mathrm{int}}(m,u)
:=
\sqrt{\mathcal W_{4,4}(m;0)\mathcal W_{4,4}(m;u)}.
\]
The regularized weights satisfy
\[
\mathcal W_{4,4}(m;0)=(1+m^{-2})^9\ge1
\]
and
\[
\begin{split}
\mathcal W_{4,4}(m;u)
&=\bigl((1-u+m^{-2})(1+u+m^{-2})\bigr)^{9/2}\\
&\ge\Delta(u)^{9/2}\ge\Delta(u)^6.
\end{split}
\]
Thus
\[
\mathcal B_{\mathrm{int}}(m,u)\ge\Delta(u)^3\ge\Delta(u)^4.
\]
The angular distance dominates the coordinate distance:
\[
|u|
=
|\cos(\arccos u)-\cos(\arccos0)|
\le|\arccos u-\arccos0|
=
\delta_{\mathrm{ang}}(0,u).
\]
Localization with exponent \(2\) therefore gives
\[
\begin{split}
|\Delta(u)^4\widetilde K_m(u)|
&\le
\frac{\Delta(u)^4C_{\mathrm{loc},2}^{4,4}m}
{\mathcal B_{\mathrm{int}}(m,u)
 (1+m\delta_{\mathrm{ang}}(0,u))^2}\\
&\le C_{\mathrm{loc},2}^{4,4}m(1+m|u|)^{-2}.
\end{split}
\]
Dividing by the normalization lower bound yields
\[
|\psi_m(u)|
\le
\frac{C_{\mathrm{loc},2}^{4,4}}{c_{\mathrm{int}}}
(1+m|u|)^{-2}.
\]
This also holds outside \([-1,1]\), where the packet vanishes.

For \(|u|<1\), symmetry of the kernel and the local angular estimate, with reference and auxiliary arguments both equal to \(u\) and fixed second argument \(0\), give, for \(v\) sufficiently close to \(u\),
\[
|\widetilde K_m(v)-\widetilde K_m(u)|
\le
\frac{C_{\mathrm{lip},1}^{4,4}m^2}
{\mathcal B_{\mathrm{int}}(m,u)
 (1+m\delta_{\mathrm{ang}}(0,u))}
|\arccos v-\arccos u|.
\]
Passing to difference-quotient limits and using
\[
\left|\frac{d}{du}\arccos u\right|
=\frac1{\sqrt{\Delta(u)}}
\]
gives
\[
|\widetilde K_m'(u)|
\le
\frac{C_{\mathrm{lip},1}^{4,4}m^2}
{\mathcal B_{\mathrm{int}}(m,u)\sqrt{\Delta(u)}}.
\]
Localization with exponent \(1\) also gives
\[
|\widetilde K_m(u)|
\le
\frac{C_{\mathrm{loc},1}^{4,4}m}
{\mathcal B_{\mathrm{int}}(m,u)}.
\]
The product rule is
\[
\psi_m'(u)
=
\frac{-8u\Delta(u)^3\widetilde K_m(u)
+\Delta(u)^4\widetilde K_m'(u)}
{\widetilde K_m(0)}.
\]
Using
\[
\frac{\Delta(u)^3}{\mathcal B_{\mathrm{int}}(m,u)}\le1,
\qquad
\frac{\Delta(u)^4}{\sqrt{\Delta(u)}}\le\Delta(u)^3,
\]
we obtain
\[
\begin{split}
|\psi_m'(u)|
&\le
\frac{8C_{\mathrm{loc},1}^{4,4}m
      +C_{\mathrm{lip},1}^{4,4}m^2}
     {\widetilde K_m(0)}\\
&\le
\frac{8C_{\mathrm{loc},1}^{4,4}
      +C_{\mathrm{lip},1}^{4,4}}{c_{\mathrm{int}}}\,m.
\end{split}
\]
Here \(m\ge1\) absorbs the term independent of \(m\) after division.

On the whole line,
\[
\psi_m(u)
=
\mathcal T(u)\frac{\widetilde K_m(u)}{\widetilde K_m(0)}.
\]
The previously established regularity of \(\mathcal T\) proves that this packet is continuously differentiable. Its derivative vanishes outside \([-1,1]\), and continuity extends the bound to the two endpoints.

Set
\[
C_{\mathrm{env,int}}
:=
\frac{C_{\mathrm{loc},2}^{4,4}
      +8C_{\mathrm{loc},1}^{4,4}
      +C_{\mathrm{lip},1}^{4,4}}{c_{\mathrm{int}}}>0.
\]
We have proved, for every real \(u\),
\[
|\psi_m(u)|\le C_{\mathrm{env,int}}(1+m|u|)^{-2},
\qquad
|\psi_m'(u)|\le C_{\mathrm{env,int}}m.
\]
% lean: CausalSmith.Stat.NoisydoseWeakdesignTransition.packetinterior_bounds_of_normalization

\item \emph{Interior symmetry, masses, and normalization at zero.}
Reindexing the binomial representation in \cref{obj:lem:jacobi-binomial} by \(i\mapsto j-i\) gives
\[
\mathsf J_j^{(4,4)}(-u)
=(-1)^j\mathsf J_j^{(4,4)}(u).
\]
At \(u=0\), every odd-degree center value is therefore zero. The surviving even-degree terms give
\[
\widetilde K_m(-u)=\widetilde K_m(u),
\qquad
\psi_m(-u)=\psi_m(u).
\]

For the exponent \(\kappa=0\), we use \(|u|^0=1\), including at \(u=0\). Define
\[
C_{\mathrm{mid,int}}:=2C_{\mathrm{env,int}}>0.
\]
The envelope bounds the supremum by \(C_{\mathrm{mid,int}}\), and the derivative is bounded by \(C_{\mathrm{mid,int}}m\). Evenness and the explicit quadratic-envelope integral give
\[
\begin{split}
\int_{-1}^1|\psi_m(u)|\,du
&=
2\int_0^1|\psi_m(u)|\,du\\
&\le
2C_{\mathrm{env,int}}\int_0^1(1+mu)^{-2}\,du
\le C_{\mathrm{mid,int}}m^{-1}.
\end{split}
\]
The pointwise inequality
\(\psi_m(u)^2\le C_{\mathrm{mid,int}}|\psi_m(u)|\) then gives
\[
\int_{-1}^1\psi_m(u)^2\,du
\le
C_{\mathrm{mid,int}}\int_{-1}^1|\psi_m(u)|\,du
\le C_{\mathrm{mid,int}}^2m^{-1}.
\]
Thus, with
\[
C_{\mathrm{fin,int}}
:=
\max\{C_{\mathrm{mid,int}},C_{\mathrm{mid,int}}^2\}>0,
\]
all four analytic bounds hold with one constant. The definition and positive normalizer also give
\[
\psi_m(u)=0\quad(|u|>1),
\qquad
\psi_m(0)=\frac{\widetilde K_m(0)}{\widetilde K_m(0)}=1.
\]
% lean: CausalSmith.Stat.NoisydoseWeakdesignTransition.packetinterior_bounds

\item \emph{Interior moments and conclusion.}
Let \(j\) be an integer with \(0\le j<m\); in particular \(j\le m\). On \([-1,1]\), the taper equals the symmetric Jacobi weight:
\[
(1-u^2)^4=(1-u)^4(1+u)^4.
\]
Expanding the finite kernel sum therefore gives
\[
\begin{split}
\int_{-1}^1u^j\psi_m(u)\,du
&=
\frac1{\widetilde K_m(0)}
\sum_{\nu=0}^{2m}
\frac{\eta(\nu/m)\mathsf J_\nu^{(4,4)}(0)}
     {H_\nu^{4,4}}\\
&\qquad\qquad{}\times
\int_{-1}^1
(1-u)^4(1+u)^4
\mathsf J_\nu^{(4,4)}(u)u^j\,du
=0.
\end{split}
\]
Indeed, an inactive coefficient is zero, while an active degree satisfies
\(\nu>m\ge j\), so its integral vanishes by Jacobi orthogonality. Finite summation and compact-interval integrability justify the expansion.

The interior branch consequently satisfies every required assertion with
\[
C:=C_{\mathrm{fin,int}},
\qquad c:=1.
\]
Together with the positive-exponent branch, this proves the result for every fixed \(\kappa\in[0,2]\).
% lean: CausalSmith.Stat.NoisydoseWeakdesignTransition.weighted_packet_construction
\end{enumerate}
\end{proof}

\begin{proof}[Proof of \cref{obj:thm:observed-lower}]
Fix public constants \(K=(p_{\min},c_{\mathrm{lo}},c_{\mathrm{hi}})\), \(\beta\in(0,1]\), \(\kappa\in[0,2]\), and \(\tau>0\), under the envelope calibration in the statement. All constants chosen below may depend on \(K,\beta,\kappa,\tau\) and the fixed packet construction.

\begin{enumerate}
\item \textbf{Packet constants.}
For \(\kappa=0\), use the symmetric interior construction in \cref{obj:def:packet-handle}; for \(\kappa>0\), use the endpoint construction. In the latter case, its Jacobi parameter satisfies
\[
b=\frac{\kappa-1}{2}>-\frac12,
\]
and the remaining Jacobi parameters equal \(r=4\). The localization estimate is supplied by \citep[Theorem 2.4, equation (2.14)]{PetrushevXu2005}, and the local increment estimate by \citep[Theorem 2.2, equation (2.4)]{https://people.math.sc.edu/pencho/Publications/kpx-06-28-06-web.pdf}. The orthogonality and positive squared-norm formula are supplied by \citep[Table 18.3.1, Jacobi row and caption]{https://dlmf.nist.gov/18.3}, and the required Gamma-ratio asymptotic by \citep[equation (5.11.12), restricted to the positive real axis]{https://dlmf.nist.gov/5.11.E12}.

Thus \cref{obj:thm:weighted-packet-construction} supplies constants
\[
C_{\mathrm{packet}}>0,\qquad c_p>0
\]
uniformly for \(m\ge4\), including positive normalizing values, evenness, continuous differentiability of the zero extension, support in \([-1,1]\), and
\[
\psi_m(0)=1,\qquad
\|\psi_m\|_\infty\le C_{\mathrm{packet}},\qquad
\|\psi_m'\|_\infty\le C_{\mathrm{packet}}m,
\]
\[
\int_{-1}^{1}|u|^\kappa|\psi_m(u)|\,du
\le C_{\mathrm{packet}}m^{-\kappa-1},
\qquad
\int_{-1}^{1}|u|^\kappa\psi_m(u)^2\,du
\le C_{\mathrm{packet}}m^{-\kappa-1},
\]
\[
\int_{-1}^{1}|u|^\kappa u^j\psi_m(u)\,du=0
\qquad(0\le j<c_pm).
\]
% lean: CausalSmith.Stat.NoisydoseWeakdesignTransition.lower_regime_witnesses

\item \textbf{Common calibration.}
Define the cancellation constant, degree threshold, intermediate multiplier, and compact radius by
\[
\gamma:=\frac{c_p}{2},\qquad
M_{\mathrm{deg}}:=\left\lceil\max\left\{4,\frac{2}{c_p}\right\}\right\rceil,
\qquad
C_{\mathrm I}:=M_{\mathrm{deg}},\qquad
\ell_{\mathrm M}:=\frac18.
\]
Then
\[
M_{\mathrm{deg}}\ge4,\qquad
c_pM_{\mathrm{deg}}\ge2,\qquad
C_{\mathrm I}\ge4,\qquad c_pC_{\mathrm I}\ge2.
\]
Choose the intermediate radius constant explicitly as
\[
a_{\mathrm I}:=
\min\left\{
1,\frac{1}{4\sqrt{C_{\mathrm I}+1}},
\sqrt{\min\left\{\frac{\gamma}{2},
\gamma\exp(-1/\gamma-1)\right\}}
\right\}.
\]
Every entry in this minimum is positive, and hence
\[
0<a_{\mathrm I}\le1,\qquad
a_{\mathrm I}\sqrt{C_{\mathrm I}+1}\le\frac14,\qquad
a_{\mathrm I}^2\le\frac{\gamma}{2},\qquad
\frac{\exp(1)a_{\mathrm I}^2}{\gamma}\le\exp(-1/\gamma).
\]
For the compact regime, set
\[
C_{\mathrm M}:=
\max\left\{
1,\frac8\gamma,\frac{2\exp(2)\ell_{\mathrm M}^2}{\gamma}
\right\}.
\]
Consequently,
\[
C_{\mathrm M}\ge1,\qquad
\frac{\gamma C_{\mathrm M}}{\exp(1)\ell_{\mathrm M}^2}
\ge2\exp(1),\qquad
\gamma C_{\mathrm M}\ge8.
\]

Define the design coefficient and three comparison constants by
\[
c_\kappa:=(\kappa+1)2^\kappa,\qquad
B_{\mathrm D}:=8c_\kappa,\qquad
B_{\mathrm I}:=4\exp(1/2)(c_\kappa C_{\mathrm{packet}})^2,\qquad
B_{\mathrm G}:=4\exp(1/2)(2c_\kappa)^2,
\]
and set
\[
B_{\mathrm{amp}}:=\max\{B_{\mathrm D},2B_{\mathrm I}\},\qquad
\varepsilon:=
\min\left\{
\frac18,\frac{1}{8C_{\mathrm{packet}}},
\sqrt{\frac{\log(1+\tau^2)}{B_{\mathrm{amp}}}}
\right\}.
\]
These constants are positive. The amplitude satisfies
\[
\varepsilon\le\frac18,\qquad
4\varepsilon\le1,\qquad
\varepsilon C_{\mathrm{packet}}\le\frac18,
\]
\[
B_{\mathrm D}\varepsilon^2\le\log(1+\tau^2),
\qquad
2B_{\mathrm I}\varepsilon^2\le\log(1+\tau^2).
\]
Finally, choose
\[
C:=\varepsilon^2\max\{B_{\mathrm G},B_{\mathrm I}\},
\qquad
c:=2\varepsilon\min\left\{
1,
\left(\frac{a_{\mathrm I}}{\sqrt{C_{\mathrm I}+1}}\right)^\beta,
\left(\frac{\ell_{\mathrm M}}{2C_{\mathrm M}}\right)^\beta
\right\}.
\]
Both are positive, and these definitions give the comparison constants required in each regime.
% lean: CausalSmith.Stat.NoisydoseWeakdesignTransition.lower_common_amplitude_constant

\item \textbf{Uniform sample cutoff.}
Use the effective exponent and sample logarithm from \cref{obj:def:frontier-rate}, and define
\[
d:=2\beta+\kappa+1>0,\qquad
L_n:=\log(\exp(1)n),\qquad h_0:=n^{-1/d},
\]
\[
n_{\mathrm D}:=\max\{2,\lceil4^d\rceil\},\qquad
n_{\log}:=\max\{2,\lceil\exp(M_{\mathrm{deg}}^2)\rceil\},
\qquad
n_0:=\max\{n_{\mathrm D},n_{\log}\}.
\]
For every \(n\ge n_0\),
\[
0<h_0\le(4^d)^{-1/d}=\frac14,\qquad
nh_0^d=1,\qquad
M_{\mathrm{deg}}^2\le L_n.
\]
The last inequality follows by taking logarithms of
\[
\exp(M_{\mathrm{deg}}^2)\le n\le\exp(1)n.
\]
For every \(\sigma\in[0,1/4]\), also
\[
1\le\log(\exp(1)+n\sigma^d)\le L_n.
\]
Indeed, \(0\le\sigma^d\le1\), while \(n\ge2\) and \(\exp(1)\ge2\) give
\[
\exp(1)\le\exp(1)+n\sigma^d
\le\exp(1)+n\le\exp(1)n.
\]
Fix henceforth an arbitrary public space from \cref{obj:def:potential-path-space}, \(n\ge n_0\), and \(\sigma\in[0,1/4]\).
% lean: CausalSmith.Stat.NoisydoseWeakdesignTransition.lower_log_sample_cutoff

\item \textbf{Witness laws and observed densities.}
The common centered-dose density is
\[
g(t):=c_\kappa|t|^\kappa\mathbf1\{|t|\le1/2\},
\qquad t\in\mathbb R.
\]
We use the convention \(|t|^0=1\), including at \(t=0\). Its normalization follows from
\[
\int_{-1/2}^{1/2}|t|^\kappa\,dt
=\frac{2^{-\kappa}}{\kappa+1},
\qquad
\int_{\mathbb R}g(t)\,dt=1.
\]
Take mutually independent seed coordinates with
\[
P(X=0)=P(X=1)=\frac12,\qquad
\mathcal L(A-a_0)(dt)=g(t)\,dt,\qquad
Z\sim N(0,1),\qquad U\sim\operatorname{Unif}[0,1].
\]
Here the prescribed center is \(a_0=1/2\). For any continuous profile
\(\mu:[0,1]\to[1/8,7/8]\), define its structural law by the measurable pushforward
\[
P_\mu:=
\mathcal L\bigl(X,A,Z,\iota(\mu,U),\iota(\mu,U)(A),U\bigr).
\]
Thus, under \(P_\mu\),
\[
Y(\cdot)=\iota(\mu,U),\qquad
Y(a)=\mathbf1\{U\le\mu(a)\},\qquad
Y=Y(A),\qquad W=A+\sigma Z.
\]
The measurability requirements follow from the embedding and joint evaluation in \cref{obj:def:potential-path-space}. Uniformity of the independent rank gives
\[
E_{P_\mu}[Y(a)\mid X=x]=\mu(a),
\qquad
\theta(P_\mu)=\sum_{x\in\{0,1\}}\frac12\mu(a_0)=\mu(a_0).
\]

For a continuous compactly supported perturbation, write
\[
\mu_\pm(a):=\frac12\pm\delta(a-a_0),\qquad
\mu_\star(a):=\frac12,\qquad
P_\pm:=P_{\mu_\pm},\qquad P_\star:=P_{\mu_\star},
\]
whenever the profiles have the displayed range. Define
\[
Q_\pm:=\mathcal L_{P_\pm}(X,W,Y),\qquad
Q_\star:=\mathcal L_{P_\star}(X,W,Y),\qquad
v_\delta(t):=g(t)\delta(t).
\]
For \(\sigma>0\), define
\[
\phi_\sigma(w):=
\frac{1}{\sqrt{2\pi}\sigma}
\exp\left(-\frac{w^2}{2\sigma^2}\right),
\qquad
p_V:=g*\phi_\sigma,\qquad f_\delta:=v_\delta*\phi_\sigma.
\]
For \(\sigma=0\), define instead
\[
p_V:=g,\qquad f_\delta:=v_\delta.
\]
Conditioning on the latent dose and integrating the uniform rank gives the densities
\[
q_\pm(x,w,y)
=
p_x\left\{\frac{p_V(w-a_0)}2
\pm(2y-1)f_\delta(w-a_0)\right\},
\qquad
q_\star(x,w,y)=\frac{p_xp_V(w-a_0)}2,
\]
relative to counting measure in \(x\in\{0,1\}\), \(y\in\{0,1\}\), and Lebesgue measure in \(w\in\mathbb R\), with \(p_x=1/2\).

The marked probabilities belong to \([0,1]\), so
\[
0\le Q_\pm\le2Q_\star,\qquad Q_\pm\ll Q_\star.
\]
For the likelihood ratios
\[
\Lambda_\pm:=\frac{dQ_\pm}{dQ_\star},
\]
this implies
\[
0\le\Lambda_\pm\le2\quad Q_\star\text{-almost surely},
\qquad
\int(\Lambda_\pm-1)^2\,dQ_\star\le1.
\]
We use
\[
\chi^2(Q_\pm,Q_\star):=
\int(\Lambda_\pm-1)^2\,dQ_\star,
\qquad
\chi^2(P_\pm,P_\star):=\chi^2(Q_\pm,Q_\star),
\]
in agreement with the theorem's observed-law convention. Summing the squared likelihood deviations over the two marks and strata, and translating \(w\) by \(a_0\), yields
\[
\chi^2(Q_\pm,Q_\star)
=
4\int_{\{p_V>0\}}\frac{f_\delta(w)^2}{p_V(w)}\,dw.
\]
At \(\sigma=0\), both marked densities vanish wherever \(g=0\), and this identity reduces to
\[
\chi^2(Q_\pm,Q_\star)
=4\int_{\mathbb R}g(t)\delta(t)^2\,dt.
\]
% lean: CausalSmith.Stat.NoisydoseWeakdesignTransition.observed_marked_chiSq_identity

\item \textbf{Gaussian moment comparison.}
Suppose \(\sigma>0\), and let the support radius and cancellation cutoff satisfy
\[
\ell>0,\qquad J\in\mathbb N_0,\qquad
\operatorname{supp}(v_\delta)\subseteq[-\ell,\ell],\qquad
\int_{\mathbb R}t^jv_\delta(t)\,dt=0\quad(0\le j<J).
\]
Define the absolute mass and moment tail by
\[
\|v_\delta\|_1:=\int_{\mathbb R}|v_\delta(t)|\,dt,\qquad
\mathcal T_{\sigma,\ell,J}:=
\sum_{j=J}^{\infty}\frac{(\ell^2/\sigma^2)^j}{j!}.
\]

First consider \(w\ge0\). Restricting the convolution defining \(p_V(w)\) to \(0\le t\le\sigma\), which lies in the latent support, gives
\[
\begin{aligned}
p_V(w)
&=\phi_\sigma(w)\int_{\mathbb R}
g(t)\exp\left(\frac{wt}{\sigma^2}-\frac{t^2}{2\sigma^2}\right)\,dt\\
&\ge
c_\kappa\exp(-1/2)\phi_\sigma(w)\int_0^\sigma t^\kappa\,dt\\
&=
2^\kappa\exp(-1/2)\sigma^{\kappa+1}\phi_\sigma(w)\\
&\ge\exp(-1/2)\sigma^{\kappa+1}\phi_\sigma(w)>0.
\end{aligned}
\]
For \(w<0\), evenness of \(g\) and \(\phi_\sigma\) gives the same bound by replacing \(w\) with \(-w\). Hence the denominator bound holds on the entire real line.

Completing the square gives, for \(t,s,w\in\mathbb R\),
\[
\frac{\phi_\sigma(w-t)\phi_\sigma(w-s)}{\phi_\sigma(w)}
=
\exp(ts/\sigma^2)\phi_\sigma(w-t-s).
\]
Its integral in \(w\) is \(\exp(ts/\sigma^2)\). Compact support bounds the absolute triple integral by
\[
\exp(\ell^2/\sigma^2)\|v_\delta\|_1^2<\infty.
\]
Absolute Fubini therefore yields
\[
\int_{\mathbb R}\frac{f_\delta(w)^2}{\phi_\sigma(w)}\,dw
=
\iint_{\mathbb R^2}
v_\delta(t)v_\delta(s)\exp(ts/\sigma^2)\,dt\,ds.
\]
For each \(j\in\mathbb N_0\), the absolute integral of the \(j\)-th exponential-series term is bounded by
\[
\iint_{\mathbb R^2}
\frac{|t^jv_\delta(t)s^jv_\delta(s)|}{\sigma^{2j}j!}\,dt\,ds
\le
\|v_\delta\|_1^2\frac{(\ell^2/\sigma^2)^j}{j!}.
\]
The sum of these bounds is finite. Termwise integration and the cancellations give
\[
\begin{aligned}
\int_{\mathbb R}\frac{f_\delta(w)^2}{\phi_\sigma(w)}\,dw
&=
\sum_{j=0}^{\infty}
\frac{\left(\int_{\mathbb R}t^jv_\delta(t)\,dt\right)^2}
{\sigma^{2j}j!}\\
&=
\sum_{j=J}^{\infty}
\frac{\left(\int_{\mathbb R}t^jv_\delta(t)\,dt\right)^2}
{\sigma^{2j}j!}\\
&\le \|v_\delta\|_1^2\mathcal T_{\sigma,\ell,J},
\end{aligned}
\]
where the last inequality uses
\[
\left|\int_{\mathbb R}t^jv_\delta(t)\,dt\right|
\le\ell^j\|v_\delta\|_1.
\]
Combining this with the observed-density identity and the denominator bound proves
\[
\chi^2(Q_\pm,Q_\star)
\le
4\exp(1/2)\sigma^{-\kappa-1}
\|v_\delta\|_1^2\mathcal T_{\sigma,\ell,J}.
\]
% lean: CausalSmith.Stat.NoisydoseWeakdesignTransition.observed_perturbation_chiSq_moment_bound

\item \textbf{Direct profiles.}
Suppose first that \(\sigma\le h_0\). Set
\[
h=\ell=h_0,\qquad m=0,\qquad J=0,
\]
and define, for all \(t\in\mathbb R\),
\[
\delta(t):=
\varepsilon h_0^\beta
\left(1-\left(\frac{t}{h_0}\right)^2\right)^2
\mathbf1\{|t|\le h_0\}.
\]
For the increment calculation, define the clipped coordinate
\[
\zeta_{\mathrm D}(t):=\min\{|t/h_0|,1\}.
\]
Then
\[
\delta(t)=\varepsilon h_0^\beta
(1-\zeta_{\mathrm D}(t)^2)^2,
\qquad
0\le\zeta_{\mathrm D}(t)\le1,
\qquad
|\zeta_{\mathrm D}(s)-\zeta_{\mathrm D}(t)|
\le\frac{|s-t|}{h_0}.
\]
Factoring differences of squares twice gives
\[
\begin{aligned}
\left|(1-\zeta_{\mathrm D}(s)^2)^2
-(1-\zeta_{\mathrm D}(t)^2)^2\right|
&\le2|\zeta_{\mathrm D}(s)^2-\zeta_{\mathrm D}(t)^2|\\
&\le4|\zeta_{\mathrm D}(s)-\zeta_{\mathrm D}(t)|\\
&\le4\frac{|s-t|}{h_0}.
\end{aligned}
\]
Thus
\[
0\le\delta(t)\le\varepsilon h_0^\beta,\qquad
|\delta(s)-\delta(t)|
\le4\varepsilon h_0^\beta\frac{|s-t|}{h_0}.
\]
If \(|s-t|\le h_0\), then
\[
\frac{|s-t|}{h_0}
\le\left(\frac{|s-t|}{h_0}\right)^\beta,
\]
so the increment is at most \(4\varepsilon|s-t|^\beta\). If \(|s-t|>h_0\), the range bound gives
\[
|\delta(s)-\delta(t)|
\le\varepsilon h_0^\beta
\le\varepsilon|s-t|^\beta.
\]
Consequently,
\[
|\delta(s)-\delta(t)|\le4\varepsilon|s-t|^\beta.
\]
The explicit formula is continuous across its support endpoints. Since \(h_0^\beta\le1\), the common calibration ensures
\[
\mu_\pm(a)\in[1/8,7/8],\qquad
|\mu_\pm(a)-\mu_\pm(a')|\le|a-a'|^\beta
\quad(a,a'\in[0,1]).
\]
At the target,
\[
|\theta(P_+)-\theta(P_-)|
=2\delta(0)=2\varepsilon h_0^\beta
=2\varepsilon\rho_{n,\sigma}
\ge c\rho_{n,\sigma},
\]
using the direct branch of \cref{obj:def:frontier-rate}.

On \([-h_0,h_0]\), the bounds on \(g\) and \(\delta\) give
\[
|v_\delta(t)|\le c_\kappa\varepsilon h_0^{\beta+\kappa}.
\]
Outside that interval \(v_\delta=0\). Integrating over its length \(2h_0\) yields
\[
\|v_\delta\|_1
\le2c_\kappa\varepsilon h_0^{\beta+\kappa+1}.
\]
Moreover,
\[
\int_{\mathbb R}g(t)\delta(t)^2\,dt
\le\varepsilon h_0^\beta\|v_\delta\|_1
\le2c_\kappa\varepsilon^2h_0^d.
\]
For positive noise, the Gaussian moment comparison with \(J=0\) therefore gives
\[
\chi^2(Q_\pm,Q_\star)
\le
B_{\mathrm G}\varepsilon^2
\sigma^{-\kappa-1}h_0^{2\beta+2\kappa+2}
\mathcal T_{\sigma,h_0,0}
\le
C\sigma^{-\kappa-1}h^{2\beta+2\kappa+2}
\mathcal T_{\sigma,\ell,J}.
\]
The moment conditions for \(J=0\) have an empty index set.
% lean: CausalSmith.Stat.NoisydoseWeakdesignTransition.direct_lower_profile_frontier_construction

\item \textbf{Direct observed testing bound.}
At \(\sigma=0\), the exact observed comparison already gives
\[
\chi^2(Q_\pm,Q_\star)
=4\int_{\mathbb R}g(t)\delta(t)^2\,dt.
\]
At \(\sigma>0\), Cauchy--Schwarz under the finite positive measure
\(g(t)\phi_\sigma(w-t)\,dt\) gives
\[
f_\delta(w)^2
\le
p_V(w)\int_{\mathbb R}
g(t)\delta(t)^2\phi_\sigma(w-t)\,dt.
\]
Dividing by the positive denominator, integrating, and using the unit Gaussian integral yields
\[
\chi^2(Q_\pm,Q_\star)
\le4\int_{\mathbb R}g(t)\delta(t)^2\,dt.
\]
Thus, for every noise value in the direct regime,
\[
n\chi^2(Q_\pm,Q_\star)
\le8c_\kappa\varepsilon^2nh_0^d
=B_{\mathrm D}\varepsilon^2
\le\log(1+\tau^2).
\]

Here is the product comparison that will also be used for the inverse regimes. Independence in \cref{obj:def:observed-experiment} and the square-integrability of the likelihood ratios imply
\[
\begin{aligned}
1+\chi^2(Q_{P_\pm}^n,Q_{P_\star}^n)
&=
E_{Q_\star^{\otimes n}}
\left[\prod_{i=1}^n\Lambda_\pm(O_i)^2\right]\\
&=
\left(E_{Q_\star}\Lambda_\pm^2\right)^n
=\{1+\chi^2(Q_\pm,Q_\star)\}^n\\
&\le\exp\{n\chi^2(Q_\pm,Q_\star)\}.
\end{aligned}
\]
Hence the sample budget implies
\[
\chi^2(Q_{P_\pm}^n,Q_{P_\star}^n)\le\tau^2.
\]
For an absolutely continuous probability law, the signed likelihood deviation has integral zero. Its positive and negative parts consequently have equal integrals, and Cauchy--Schwarz gives
\[
\operatorname{TV}(Q_{P_\pm}^n,Q_{P_\star}^n)
=
\frac12E_{Q_{P_\star}^n}
\left|\prod_{i=1}^n\Lambda_\pm(O_i)-1\right|
\le\frac12\sqrt{\chi^2(Q_{P_\pm}^n,Q_{P_\star}^n)}
\le\frac\tau2.
\]
The triangle inequality now proves
\[
\operatorname{TV}(Q_{P_+}^n,Q_{P_-}^n)\le\tau.
\]
% lean: CausalSmith.Stat.NoisydoseWeakdesignTransition.direct_profile_pair_observed_tv_bound

\item \textbf{Inverse profiles and tail estimates.}
For the remaining cases, \(\sigma>h_0>0\). Given
\[
m\in\mathbb N,\qquad m\ge M_{\mathrm{deg}},\qquad
0<\ell\le\frac14,
\]
define
\[
h:=\frac{\ell}{m},\qquad
J:=\lfloor c_pm\rfloor,\qquad
\delta(t):=\varepsilon h^\beta\psi_m(t/\ell)
\quad(t\in\mathbb R).
\]
Then \(0<h\le1/4\). Since \(c_pm\ge2\) and
\(\lfloor c_pm\rfloor>c_pm-1\),
\[
J\ge\gamma m,\qquad J\ge1.
\]
The packet bounds imply continuity and support in \([-\ell,\ell]\), together with
\[
|\delta(t)|\le\varepsilon C_{\mathrm{packet}}h^\beta.
\]
The derivative bound and the mean-value inequality give
\[
|\delta(s)-\delta(t)|
\le\varepsilon C_{\mathrm{packet}}h^\beta\frac{|s-t|}{h}.
\]
If \(|s-t|\le h\), replacing \(|s-t|/h\) by its \(\beta\)-th power proves
\[
|\delta(s)-\delta(t)|
\le\varepsilon C_{\mathrm{packet}}|s-t|^\beta.
\]
If \(|s-t|>h\), the supremum bound gives
\[
|\delta(s)-\delta(t)|
\le2\varepsilon C_{\mathrm{packet}}h^\beta
\le2\varepsilon C_{\mathrm{packet}}|s-t|^\beta.
\]
Therefore
\[
|\delta(s)-\delta(t)|
\le2\varepsilon C_{\mathrm{packet}}|s-t|^\beta.
\]
The common amplitude calibration implies
\[
\mu_\pm(a)\in[1/8,7/8],\qquad
|\mu_\pm(a)-\mu_\pm(a')|\le|a-a'|^\beta.
\]
These arguments include the Lipschitz endpoint \(\beta=1\).

For \(j<J\), we have \(j<c_pm\). Changing variables \(t=\ell u\), using \(\ell\le1/4\), and applying the packet cancellations gives
\[
\int_{\mathbb R}t^jv_\delta(t)\,dt
=
c_\kappa\varepsilon h^\beta\ell^{\kappa+j+1}
\int_{-1}^{1}|u|^\kappa u^j\psi_m(u)\,du
=0.
\]
The weighted absolute-mass bound similarly gives
\[
\begin{aligned}
\|v_\delta\|_1
&=
c_\kappa\varepsilon h^\beta\ell^{\kappa+1}
\int_{-1}^{1}|u|^\kappa|\psi_m(u)|\,du\\
&\le
c_\kappa\varepsilon C_{\mathrm{packet}}
h^\beta\ell^{\kappa+1}m^{-\kappa-1}\\
&=
c_\kappa\varepsilon C_{\mathrm{packet}}h^{\beta+\kappa+1}.
\end{aligned}
\]
Thus the Gaussian moment comparison proves
\[
\chi^2(Q_\pm,Q_\star)
\le
B_{\mathrm I}\varepsilon^2
\sigma^{-\kappa-1}h^{2\beta+2\kappa+2}
\mathcal T_{\sigma,\ell,J}
\le
C\sigma^{-\kappa-1}h^{2\beta+2\kappa+2}
\mathcal T_{\sigma,\ell,J}.
\]
Normalization at zero gives the exact separation
\[
|\theta(P_+)-\theta(P_-)|
=2\varepsilon h^\beta.
\]

To control the sample budget, define
\[
\lambda:=\frac{\ell^2}{\sigma^2}>0.
\]
Whenever \(J\ge1\) and \(J\ge2\lambda\), consecutive factorial-series terms have ratio at most \(1/2\). Induction therefore gives
\[
\frac{\lambda^{J+k}}{(J+k)!}
\le\frac{\lambda^J}{J!}\,2^{-k}
\quad(k\in\mathbb N_0),
\qquad
\mathcal T_{\sigma,\ell,J}
\le2\frac{\lambda^J}{J!}.
\]
The nonnegative exponential series gives
\[
\frac{J^J}{J!}
\le\sum_{k=0}^\infty\frac{J^k}{k!}
=\exp(J).
\]
Multiplication by \((\lambda/J)^J\) yields
\[
\frac{\lambda^J}{J!}
\le\left(\frac{\exp(1)\lambda}{J}\right)^J,
\qquad
\mathcal T_{\sigma,\ell,J}
\le2\left(\frac{\exp(1)\lambda}{J}\right)^J.
\]
Finally, whenever \(h\le\sigma\),
\[
\sigma^{-\kappa-1}h^{2\beta+2\kappa+2}
=
h^d(h/\sigma)^{\kappa+1}
\le h^d.
\]
% lean: CausalSmith.Stat.NoisydoseWeakdesignTransition.inverse_lower_profile_chiSq_construction

\item \textbf{Intermediate regime.}
Suppose
\[
h_0<\sigma\le L_n^{-1/2}.
\]
Define the effective sample quantity and its logarithm by
\[
N_{\mathrm I}:=n\sigma^d,\qquad
S_{\mathrm I}:=\log(\exp(1)+N_{\mathrm I}),
\]
and choose
\[
m:=\lceil C_{\mathrm I}S_{\mathrm I}\rceil,\qquad
\ell:=a_{\mathrm I}\sigma\sqrt m,\qquad
h:=\frac{\ell}{m}=\frac{a_{\mathrm I}\sigma}{\sqrt m},
\qquad
J:=\lfloor c_pm\rfloor.
\]
The logarithmic window gives \(1\le S_{\mathrm I}\le L_n\). The ceiling bounds imply
\[
C_{\mathrm I}S_{\mathrm I}\le m
<C_{\mathrm I}S_{\mathrm I}+1
\le(C_{\mathrm I}+1)S_{\mathrm I}.
\]
In particular \(m\ge M_{\mathrm{deg}}\), so \(J\ge\gamma m\) and \(J\ge1\). Moreover,
\[
\begin{aligned}
0<\ell
&\le a_{\mathrm I}\sigma
\sqrt{C_{\mathrm I}+1}\sqrt{L_n}\\
&\le a_{\mathrm I}\sqrt{C_{\mathrm I}+1}
\le\frac14,
\end{aligned}
\qquad
0<h\le\ell\le\frac14,\qquad h\le\sigma.
\]
Thus the inverse profiles just constructed are legal. Their bandwidth satisfies
\[
h=\frac{a_{\mathrm I}\sigma}{\sqrt m}
\ge
\frac{a_{\mathrm I}}{\sqrt{C_{\mathrm I}+1}}
\frac{\sigma}{\sqrt{S_{\mathrm I}}},
\]
and therefore
\[
|\theta(P_+)-\theta(P_-)|
\ge
2\varepsilon
\left(\frac{a_{\mathrm I}}{\sqrt{C_{\mathrm I}+1}}\right)^\beta
\left(\frac{\sigma}{\sqrt{S_{\mathrm I}}}\right)^\beta
\ge c\rho_{n,\sigma},
\]
using the intermediate branch of \cref{obj:def:frontier-rate}.

For the Gaussian tail parameter,
\[
\lambda=\frac{\ell^2}{\sigma^2}=a_{\mathrm I}^2m,
\qquad
2\lambda\le\gamma m\le J,
\]
and the radius calibration gives
\[
\frac{\exp(1)\lambda}{J}
\le\frac{\exp(1)a_{\mathrm I}^2}{\gamma}
\le\exp(-1/\gamma).
\]
Consequently,
\[
\mathcal T_{\sigma,\ell,J}
\le2\exp(-J/\gamma)\le2\exp(-m).
\]
Since \(m\ge S_{\mathrm I}\),
\[
N_{\mathrm I}\exp(-m)
\le
N_{\mathrm I}\exp(-S_{\mathrm I})
=
\frac{N_{\mathrm I}}{\exp(1)+N_{\mathrm I}}
\le1.
\]
Using \(h\le\sigma\) and \(d>0\), we obtain the complete sample-factor bound
\[
\begin{aligned}
n\sigma^{-\kappa-1}h^{2\beta+2\kappa+2}
\mathcal T_{\sigma,\ell,J}
&\le2n\sigma^d\exp(-m)\\
&=2N_{\mathrm I}\exp(-m)\le2.
\end{aligned}
\]
Hence
\[
n\chi^2(Q_\pm,Q_\star)
\le2B_{\mathrm I}\varepsilon^2
\le\log(1+\tau^2).
\]
The product comparison established above proves
\[
\operatorname{TV}(Q_{P_+}^n,Q_{P_-}^n)\le\tau.
\]
% lean: CausalSmith.Stat.NoisydoseWeakdesignTransition.intermediate_lower_profile_sample_budget

\item \textbf{Compact regime.}
The remaining case satisfies
\[
\sigma>h_0,\qquad \sigma>L_n^{-1/2}.
\]
Define
\[
\tau_{\mathrm M}:=\sigma^2L_n>1,\qquad
B_{\mathrm M}:=\log(\exp(1)+\tau_{\mathrm M}),
\]
and choose
\[
m:=\left\lceil\frac{C_{\mathrm M}L_n}{B_{\mathrm M}}\right\rceil,
\qquad
\ell:=\ell_{\mathrm M},\qquad
h:=\frac{\ell_{\mathrm M}}m,\qquad
J:=\lfloor c_pm\rfloor.
\]

We first establish the logarithmic bounds used for this choice. Since
\(\tau_{\mathrm M}\ge1\),
\[
\exp(1)+\tau_{\mathrm M}
\le(\exp(1)+1)\tau_{\mathrm M}.
\]
Also \(\exp(1)+1\le\exp(2)\), so
\[
B_{\mathrm M}
\le\log\tau_{\mathrm M}+\log(\exp(1)+1)
\le\log\tau_{\mathrm M}+2.
\]
Applying \(\log x\le x-1\) to \(x=\sqrt{\tau_{\mathrm M}}\) gives
\[
\log\tau_{\mathrm M}\le2\sqrt{\tau_{\mathrm M}}-2,
\qquad
B_{\mathrm M}\le2\sqrt{\tau_{\mathrm M}}.
\]
Because \(\sigma\le1/4\),
\[
0<B_{\mathrm M}
\le2\sigma\sqrt{L_n}
\le\frac{\sqrt{L_n}}2,
\qquad
\frac{L_n}{B_{\mathrm M}}\ge2\sqrt{L_n}\ge1.
\]
The ceiling bounds consequently yield
\[
\frac{C_{\mathrm M}L_n}{B_{\mathrm M}}
\le m\le\frac{2C_{\mathrm M}L_n}{B_{\mathrm M}}.
\]
Furthermore,
\[
m\ge\frac{L_n}{B_{\mathrm M}}
\ge2\sqrt{L_n}\ge M_{\mathrm{deg}},
\]
where \(M_{\mathrm{deg}}^2\le L_n\) was enforced by the sample cutoff. Thus
\[
J\ge\gamma m\ge
\frac{\gamma C_{\mathrm M}L_n}{B_{\mathrm M}},
\qquad J\ge1.
\]
The bandwidth is positive and at most \(\ell_{\mathrm M}\le1/4\). Since
\(\sigma\sqrt{L_n}\ge1\),
\[
\sigma m\ge2\sigma\sqrt{L_n}\ge2\ge\ell_{\mathrm M},
\qquad h\le\sigma.
\]
The upper degree bound gives
\[
h=\frac{\ell_{\mathrm M}}m
\ge
\frac{\ell_{\mathrm M}}{2C_{\mathrm M}}\frac{B_{\mathrm M}}{L_n}.
\]
The inverse-profile separation therefore satisfies
\[
|\theta(P_+)-\theta(P_-)|
\ge
2\varepsilon
\left(\frac{\ell_{\mathrm M}}{2C_{\mathrm M}}\right)^\beta
\left(\frac{B_{\mathrm M}}{L_n}\right)^\beta
\ge c\rho_{n,\sigma},
\]
using the compact branch of \cref{obj:def:frontier-rate}.

For the tail calculation, define
\[
\lambda:=\frac{\ell_{\mathrm M}^2}{\sigma^2}
=\frac{\ell_{\mathrm M}^2L_n}{\tau_{\mathrm M}},
\qquad
D_{\mathrm M}:=
\frac{\gamma C_{\mathrm M}}
{\exp(1)\ell_{\mathrm M}^2}\ge2\exp(1).
\]
The cancellation order and the bound \(B_{\mathrm M}\le2\sqrt{\tau_{\mathrm M}}\) imply
\[
\frac{J}{\exp(1)\lambda}
\ge
D_{\mathrm M}\frac{\tau_{\mathrm M}}{B_{\mathrm M}}
\ge\exp(1)\sqrt{\tau_{\mathrm M}}
\ge2.
\]
In particular \(J\ge2\lambda\). Taking logarithms, and using
\(\log(\exp(1)+1)\le2\), gives
\[
\log\frac{J}{\exp(1)\lambda}
\ge1+\frac12\log\tau_{\mathrm M}
\ge\frac{B_{\mathrm M}}2.
\]
Hence
\[
J\log\frac{J}{\exp(1)\lambda}
\ge
\frac{\gamma C_{\mathrm M}L_n}{B_{\mathrm M}}
\frac{B_{\mathrm M}}2
=
\frac{\gamma C_{\mathrm M}L_n}{2}
\ge4L_n.
\]
The factorial-tail estimate now gives
\[
\mathcal T_{\sigma,\ell_{\mathrm M},J}
\le
2\exp\left(-J\log\frac{J}{\exp(1)\lambda}\right)
\le2\exp(-4L_n).
\]
Using \(h\le\sigma\) for the first inequality and \(0<h\le1\) for the second, the Gaussian prefactor satisfies
\[
\sigma^{-\kappa-1}h^{2\beta+2\kappa+2}\le h^d\le1.
\]
Also \(L_n\ge0\), so
\[
n\exp(-4L_n)
\le n\exp(-L_n)
=\exp(-1)\le1.
\]
It follows that
\[
n\sigma^{-\kappa-1}h^{2\beta+2\kappa+2}
\mathcal T_{\sigma,\ell_{\mathrm M},J}\le2,
\]
and therefore
\[
n\chi^2(Q_\pm,Q_\star)
\le2B_{\mathrm I}\varepsilon^2
\le\log(1+\tau^2).
\]
The same product comparison proves
\[
\operatorname{TV}(Q_{P_+}^n,Q_{P_-}^n)\le\tau.
\]
% lean: CausalSmith.Stat.NoisydoseWeakdesignTransition.compact_lower_profile_sample_budget

\item \textbf{Structural membership and assembly.}
Each regime has supplied continuous profiles satisfying
\[
\mu_\pm(a)\in[1/8,7/8],\qquad
|\mu_\pm(a)-\mu_\pm(a')|\le|a-a'|^\beta.
\]
The reference profile satisfies the same conditions, since
\[
\mu_\star(a)=\frac12,\qquad
|\mu_\star(a)-\mu_\star(a')|=0\le|a-a'|^\beta.
\]
We now verify membership for all three pushforward laws.

The common seed gives
\[
p_0=p_1=\frac12,\qquad
A=a_0+t\in[0,1]\quad\text{almost surely}.
\]
The common conditional centered-dose density has coefficient
\(c_\kappa=(\kappa+1)2^\kappa\). The assumed calibration therefore gives
\[
c_{\mathrm{lo}}|t|^\kappa\le g(t)
\le c_{\mathrm{hi}}|t|^\kappa
\qquad(t\in[-1/2,1/2]).
\]
Because the class constants satisfy \(p_{\min}\le1/2\), the equal stratum probabilities obey \(p_{\min}\le p_x=1/2\le1-p_{\min}\) for both strata. These facts establish \cref{obj:ass:stratum-overlap,obj:ass:latent-support,obj:ass:weak-design}.

For every one of the three profiles, the conditional-mean identity derived from the uniform seed is
\[
E_{P_\mu}[Y(a)\mid X=x]=\mu(a).
\]
The established profile bounds therefore give
\cref{obj:ass:holder-mean,obj:ass:mean-range}. The chosen Gaussian marginal gives \cref{obj:ass:gaussian-channel}. Independence of the Gaussian seed from the other seed coordinates is preserved under the measurable schedule map, yielding
\[
Z\perp(X,A,Y(\cdot)),
\]
as required by \cref{obj:ass:error-independence}.

The product seed also gives
\[
\mathcal L(U\mid X=x)=\operatorname{Unif}[0,1],
\qquad U\perp A\mid X,
\]
establishing \cref{obj:ass:rank-uniform,obj:ass:rank-treatment-independence}. More explicitly, for measurable schedule and dose events,
\[
\begin{aligned}
&P_\mu\{Y(\cdot)\in\mathcal B_{\mathrm{sched}},
A\in\mathcal B_{\mathrm{dose}}\mid X=x\}\\
&\qquad=
P\{\iota(\mu,U)\in\mathcal B_{\mathrm{sched}}\}
P\{A\in\mathcal B_{\mathrm{dose}}\}.
\end{aligned}
\]
Here the events have domains
\[
\mathcal B_{\mathrm{sched}}\in\mathcal E,\qquad
\mathcal B_{\mathrm{dose}}\in\mathcal B([0,1]).
\]
This factorization proves \cref{obj:ass:conditional-unconfoundedness} for the entire random schedule.

The measurable pushforward construction carries a single full-measure event on which
\[
Y(a)=\mathbf1\{U\le\mu(a)\}
=\mathbf1\{U\le\mu_X(a)\}
\qquad\text{for every }a\in[0,1].
\]
Thus \cref{obj:ass:structural-threshold,obj:ass:binary-potential-outcomes} hold. The same identity gives \(Y(a)\in[0,1]\), and the construction gives \(Y=Y(A)\); these establish \cref{obj:ass:bounded-potential-outcomes,obj:ass:consistency}. All ten conditions in \cref{obj:def:model-class} are therefore satisfied by \(P_+,P_-,P_\star\).

For each law, take the observed product measure
\[
Q_P^n:=\bigotimes_{i=1}^n\mathcal L_P(X,W,Y).
\]
This is the experiment in \cref{obj:def:observed-experiment} and satisfies \cref{obj:ass:iid}. Finally, projection of the seed pushforward onto \((X,A)\) removes the profile-dependent coordinates, so
\[
\mathcal L_{P_+}(X,A)
=\mathcal L_{P_-}(X,A)
=\mathcal L_{P_\star}(X,A).
\]
The direct branch includes \(\sigma=0\) and \(\sigma=h_0\); among the remaining noise values, the intermediate branch includes \(\sigma=L_n^{-1/2}\), and the compact branch covers the strict upper side. Each branch has supplied the required separation, observed product-TV bound, and, for positive noise, the displayed Gaussian moment-series comparison with the same constant \(C\). The three laws use the same seed and the same regime-specific tuning data, as asserted.
% lean: CausalSmith.Stat.NoisydoseWeakdesignTransition.observed_lower
\end{enumerate}
\end{proof}
